\IfFileExists{stacks-project.cls}{%
\documentclass{stacks-project}
}{%
\documentclass{amsart}
}
\usepackage{amssymb}

% For dealing with references we use the comment environment
\usepackage{verbatim}
\usepackage{fancyvrb}
\usepackage{pmboxdraw}
\newenvironment{reference}{\comment}{\endcomment}
%\newenvironment{reference}{}{}
\newenvironment{slogan}{\comment}{\endcomment}
\newenvironment{history}{\comment}{\endcomment}

% For source figures retained as assets
\usepackage{graphicx}

% For commutative diagrams we use Xy-pic
\usepackage[all]{xy}

% We use 2cell for 2-commutative diagrams.
\xyoption{2cell}
\UseAllTwocells

% We use multicol for the list of chapters between chapters
\usepackage{multicol}

% This is generally recommended for better output
\usepackage{lmodern}
\usepackage[T1]{fontenc}
\usepackage{textalpha}

% For cross-file-references
% We prefer the snapshot tmp/xref/#2.aux,
% which latexmkrc replaces atomically, so
% that documents compiled simultaneously
% never read a half-written .aux file.
% The last argument keeps links to #2.pdf.
\usepackage{xr-hyper}
\newcommand{\maybeexternaldocument}[2]{%
  \IfFileExists{#2.aux}{%
    \externaldocument[#1]{#2}%
  }{%
    \PackageWarning{stack}{Missing #2.aux; external references with prefix #1 unavailable}%
  }%
}
\let\externaldocumentorig\externaldocument
\renewcommand{\externaldocument}[2][]{%
  \IfFileExists{tmp/xref/#2.aux}{%
    \externaldocumentorig[#1]{tmp/xref/#2}[#2.pdf]%
  }{%
  \IfFileExists{#2.aux}{%
    \externaldocumentorig[#1]{#2}%
  }{%
    \PackageWarning{stack}{Missing #2.aux; external references with prefix #1 unavailable}%
  }%
  }%
}

% Package for hypertext links:
\usepackage{hyperref}

% For any local file, say "hello.tex" you want to link to please
% use \externaldocument[hello-]{hello}
\externaldocument[introduction-]{introduction}
\externaldocument[conventions-]{conventions}
\externaldocument[sets-]{sets}
\externaldocument[categories-]{categories}
\externaldocument[topology-]{topology}
\externaldocument[sheaves-]{sheaves}
\externaldocument[sites-]{sites}
\externaldocument[stacks-]{stacks}
\externaldocument[fields-]{fields}
\externaldocument[algebra-]{algebra}
\externaldocument[brauer-]{brauer}
\externaldocument[homology-]{homology}
\externaldocument[derived-]{derived}
\externaldocument[simplicial-]{simplicial}
\externaldocument[more-algebra-]{more-algebra}
\externaldocument[smoothing-]{smoothing}
\externaldocument[modules-]{modules}
\externaldocument[sites-modules-]{sites-modules}
\externaldocument[injectives-]{injectives}
\externaldocument[cohomology-]{cohomology}
\externaldocument[sites-cohomology-]{sites-cohomology}
\externaldocument[dga-]{dga}
\externaldocument[dpa-]{dpa}
\externaldocument[sdga-]{sdga}
\externaldocument[hypercovering-]{hypercovering}
\externaldocument[schemes-]{schemes}
\externaldocument[constructions-]{constructions}
\externaldocument[properties-]{properties}
\externaldocument[morphisms-]{morphisms}
\externaldocument[coherent-]{coherent}
\externaldocument[divisors-]{divisors}
\externaldocument[limits-]{limits}
\externaldocument[varieties-]{varieties}
\externaldocument[topologies-]{topologies}
\externaldocument[descent-]{descent}
\externaldocument[perfect-]{perfect}
\externaldocument[more-morphisms-]{more-morphisms}
\externaldocument[flat-]{flat}
\externaldocument[groupoids-]{groupoids}
\externaldocument[more-groupoids-]{more-groupoids}
\externaldocument[etale-]{etale}
\externaldocument[chow-]{chow}
\externaldocument[intersection-]{intersection}
\externaldocument[pic-]{pic}
\externaldocument[weil-]{weil}
\externaldocument[adequate-]{adequate}
\externaldocument[dualizing-]{dualizing}
\externaldocument[duality-]{duality}
\externaldocument[discriminant-]{discriminant}
\externaldocument[derham-]{derham}
\externaldocument[local-cohomology-]{local-cohomology}
\externaldocument[algebraization-]{algebraization}
\externaldocument[curves-]{curves}
\externaldocument[resolve-]{resolve}
\externaldocument[models-]{models}
\externaldocument[functors-]{functors}
\externaldocument[equiv-]{equiv}
\externaldocument[pione-]{pione}
\externaldocument[etale-cohomology-]{etale-cohomology}
\externaldocument[proetale-]{proetale}
\externaldocument[relative-cycles-]{relative-cycles}
\externaldocument[more-etale-]{more-etale}
\externaldocument[trace-]{trace}
\externaldocument[crystalline-]{crystalline}
\externaldocument[spaces-]{spaces}
\externaldocument[spaces-properties-]{spaces-properties}
\externaldocument[spaces-morphisms-]{spaces-morphisms}
\externaldocument[decent-spaces-]{decent-spaces}
\externaldocument[spaces-cohomology-]{spaces-cohomology}
\externaldocument[spaces-limits-]{spaces-limits}
\externaldocument[spaces-divisors-]{spaces-divisors}
\externaldocument[spaces-over-fields-]{spaces-over-fields}
\externaldocument[spaces-topologies-]{spaces-topologies}
\externaldocument[spaces-descent-]{spaces-descent}
\externaldocument[spaces-perfect-]{spaces-perfect}
\externaldocument[spaces-more-morphisms-]{spaces-more-morphisms}
\externaldocument[spaces-flat-]{spaces-flat}
\externaldocument[spaces-groupoids-]{spaces-groupoids}
\externaldocument[spaces-more-groupoids-]{spaces-more-groupoids}
\externaldocument[bootstrap-]{bootstrap}
\externaldocument[spaces-pushouts-]{spaces-pushouts}
\externaldocument[spaces-chow-]{spaces-chow}
\externaldocument[groupoids-quotients-]{groupoids-quotients}
\externaldocument[spaces-more-cohomology-]{spaces-more-cohomology}
\externaldocument[spaces-simplicial-]{spaces-simplicial}
\externaldocument[spaces-duality-]{spaces-duality}
\externaldocument[formal-spaces-]{formal-spaces}
\externaldocument[restricted-]{restricted}
\externaldocument[spaces-resolve-]{spaces-resolve}
\externaldocument[formal-defos-]{formal-defos}
\externaldocument[defos-]{defos}
\externaldocument[cotangent-]{cotangent}
\externaldocument[examples-defos-]{examples-defos}
\externaldocument[algebraic-]{algebraic}
\externaldocument[examples-stacks-]{examples-stacks}
\externaldocument[stacks-sheaves-]{stacks-sheaves}
\externaldocument[criteria-]{criteria}
\externaldocument[artin-]{artin}
\externaldocument[quot-]{quot}
\externaldocument[stacks-properties-]{stacks-properties}
\externaldocument[stacks-morphisms-]{stacks-morphisms}
\externaldocument[stacks-limits-]{stacks-limits}
\externaldocument[stacks-cohomology-]{stacks-cohomology}
\externaldocument[stacks-perfect-]{stacks-perfect}
\externaldocument[stacks-introduction-]{stacks-introduction}
\externaldocument[stacks-more-morphisms-]{stacks-more-morphisms}
\externaldocument[stacks-geometry-]{stacks-geometry}
\externaldocument[moduli-]{moduli}
\externaldocument[moduli-curves-]{moduli-curves}
\externaldocument[examples-]{examples}
\externaldocument[exercises-]{exercises}
\externaldocument[guide-]{guide}
\externaldocument[desirables-]{desirables}
\externaldocument[coding-]{coding}
\externaldocument[obsolete-]{obsolete}
\externaldocument[fdl-]{fdl}
\externaldocument[index-]{index}

% Heap Project documents
\externaldocument[k3-]{k3}
\externaldocument[torelli-]{torelli}
\externaldocument[linear-algebra-]{linear-algebra}
\externaldocument[groups-]{groups}
\externaldocument[complex-analysis-]{complex-analysis}
\externaldocument[functional-analysis-]{functional-analysis}
\externaldocument[categories-]{categories}
\externaldocument[commutative-algebra-]{commutative-algebra}
\externaldocument[ringed-spaces-]{ringed-spaces}
\externaldocument[homological-algebra-]{homological-algebra}
\externaldocument[lie-algebras-]{lie-algebras}
\externaldocument[representation-theory-]{representation-theory}
\externaldocument[smooth-manifolds-]{smooth-manifolds}
\externaldocument[differential-topology-]{differential-topology}
\externaldocument[algebraic-topology-]{algebraic-topology}
\externaldocument[differential-geometry-]{differential-geometry}
\externaldocument[algebraic-geometry-]{algebraic-geometry}
\externaldocument[symplectic-geometry-]{symplectic-geometry}
\externaldocument[complex-geometry-]{complex-geometry}
\externaldocument[hodge-theory-]{hodge-theory}
\externaldocument[%
variations-of-hodge-structures-]{%
variations-of-hodge-structures}
\externaldocument[%
homological-mirror-symmetry-]{%
homological-mirror-symmetry}
\maybeexternaldocument{deformations-}{deformations}
\externaldocument[vertex-algebras-]{vertex-algebras}
\externaldocument[springer-]{springer}
\externaldocument[infty-categories-]{infty-categories}
\externaldocument[seminars-differential-geometry-]{%
seminars-differential-geometry}
\externaldocument[seminars-complex-geometry-]{%
seminars-complex-geometry}
\externaldocument[seminars-algebraic-geometry-]{%
seminars-algebraic-geometry}
\externaldocument[seminars-analysis-pde-]{%
seminars-analysis-pde}
\externaldocument[surfaces-]{surfaces}
\externaldocument[birational-maps-commutative-algebra-]{birational-maps-commutative-algebra}
\externaldocument[cremona-transformations-]{cremona-transformations}
\externaldocument[derived-categories-of-sheaves-]{derived-categories-of-sheaves}
\externaldocument[geometric-invariant-theory-]{geometric-invariant-theory}
\externaldocument[geometric-stability-conditions-and-group-actions-]{geometric-stability-conditions-and-group-actions}
\externaldocument[moduli-spaces-of-sheaves-]{moduli-spaces-of-sheaves}
\externaldocument[bridgeland-stability-general-theory-]{bridgeland-stability-general-theory}
\externaldocument[hilbert-schemes-]{hilbert-schemes}
\externaldocument[noncommutative-abelian-surfaces-and-kummer-type-hyperkahler-varieties-]{noncommutative-abelian-surfaces-and-kummer-type-hyperkahler-varieties}
\externaldocument[mumford-tate-groups-in-hodge-theory-]{mumford-tate-groups-in-hodge-theory}
\externaldocument[hodge-birational-atoms-2026-]{hodge-birational-atoms-2026}
\externaldocument[xxii-egd-2026-]{%
xxii-egd-2026}
\externaldocument[pde-]{pde}
\externaldocument[probability-]{probability}

\externaldocument[linguistics-]{linguistics}

\externaldocument[%
16th-alga-meeting-2026-]{%
16th-alga-meeting-2026}

\externaldocument[%
iv-jornadas-lefschetz-2026-]{%
iv-jornadas-lefschetz-2026}

\externaldocument[reading-the-masters-]{%
reading-the-masters}

\externaldocument[real-analysis-]{%
real-analysis}

\externaldocument[optimal-transport-]{%
optimal-transport}

\externaldocument[history-]{history}

\externaldocument[seiberg-witten-]{seiberg-witten}

\externaldocument[history-philosophy-biology-]{history-philosophy-biology}

\externaldocument[classical-mechanics-]{classical-mechanics}
\externaldocument[electromagnetism-]{electromagnetism}
\externaldocument[waves-and-optics-]{waves-and-optics}
\externaldocument[wif-]{waves-in-fluids-iv}
\externaldocument[thermal-physics-]{thermal-physics}
\externaldocument[quantum-mechanics-]{quantum-mechanics}
\externaldocument[special-relativity-]{special-relativity}
\externaldocument[physics-]{physics}

\externaldocument[chemistry-]{chemistry}

\externaldocument[molecular-cell-biology-]{molecular-cell-biology}

\externaldocument[microbiology-]{microbiology}

\externaldocument[protistology-]{protistology}

\externaldocument[mycology-]{mycology}

\externaldocument[botany-]{botany}

\externaldocument[zoology-]{zoology}

\externaldocument[animal-histology-]{animal-histology}

\externaldocument[animal-physiology-]{animal-physiology}

\externaldocument[systematics-biogeography-]{systematics-biogeography}

\externaldocument[ecology-natural-resources-]{ecology-natural-resources}

\externaldocument[earth-sciences-]{earth-sciences}

\externaldocument[biostatistics-]{biostatistics}

\externaldocument[mathematical-systems-biology-]{mathematical-systems-biology}

\externaldocument[cimpa-el-salvador-2024-]{cimpa-el-salvador-2024}

\externaldocument[anthropology-]{anthropology}

\externaldocument[yang-mills-theory-]{yang-mills-theory}

\externaldocument[seminars-algebra-]{seminars-algebra}

\externaldocument[seminars-physics-]{seminars-physics}

\externaldocument[seminars-computing-]{seminars-computing}

\externaldocument[oktobermat-2026-]{oktobermat-2026}

\externaldocument[saet-2026-]{saet-2026}

\externaldocument[seminars-symplectic-geometry-]{seminars-symplectic-geometry}

\externaldocument[seminars-dynamics-]{seminars-dynamics}

\externaldocument[oktobermat-2024-]{oktobermat-2024}

\externaldocument[quantum-groups-]{quantum-groups}

\externaldocument[legal-2027-]{legal-2027}

\externaldocument[derived-categories-]{derived-categories}

% Theorem environments.
%
\theoremstyle{plain}
\newtheorem{theorem}[subsection]{Theorem}
\newtheorem{proposition}[subsection]{Proposition}
\newtheorem{lemma}[subsection]{Lemma}
\newtheorem{corollary}[subsection]{Corollary}

\theoremstyle{definition}
\newtheorem{definition}[subsection]{Definition}
\newtheorem{example}[subsection]{Example}
\newtheorem{exercise}[subsection]{Exercise}
\newtheorem{situation}[subsection]{Situation}

\theoremstyle{remark}
\newtheorem{remark}[subsection]{Remark}
\newtheorem{remarks}[subsection]{Remarks}

\numberwithin{equation}{subsection}

% Macros
%
\def\lim{\mathop{\mathrm{lim}}\nolimits}
\def\colim{\mathop{\mathrm{colim}}\nolimits}
\def\Spec{\mathop{\mathrm{Spec}}}
\def\Hom{\mathop{\mathrm{Hom}}\nolimits}
\def\Ext{\mathop{\mathrm{Ext}}\nolimits}
\def\SheafHom{\mathop{\mathcal{H}\!\mathit{om}}\nolimits}
\def\SheafExt{\mathop{\mathcal{E}\!\mathit{xt}}\nolimits}
\def\Sch{\mathit{Sch}}
\def\Mor{\mathop{\mathrm{Mor}}\nolimits}
\def\Ob{\mathop{\mathrm{Ob}}\nolimits}
\def\Sh{\mathop{\mathit{Sh}}\nolimits}
\def\NL{\mathop{N\!L}\nolimits}
\def\CH{\mathop{\mathrm{CH}}\nolimits}
\def\proetale{{pro\text{-}\acute{e}tale}}
\def\etale{{\acute{e}tale}}
\def\QCoh{\mathit{QCoh}}
\def\Ker{\mathop{\mathrm{Ker}}}
\def\Im{\mathop{\mathrm{Im}}}
\def\Coker{\mathop{\mathrm{Coker}}}
\def\Coim{\mathop{\mathrm{Coim}}}

% Boxtimes
%
\DeclareMathSymbol{\boxtimes}{\mathbin}{AMSa}{"02}

%
% Macros for moduli stacks/spaces
%
\def\QCohstack{\mathcal{QC}\!\mathit{oh}}
\def\Cohstack{\mathcal{C}\!\mathit{oh}}
\def\Spacesstack{\mathcal{S}\!\mathit{paces}}
\def\Quotfunctor{\mathrm{Quot}}
\def\Hilbfunctor{\mathrm{Hilb}}
\def\Curvesstack{\mathcal{C}\!\mathit{urves}}
\def\Polarizedstack{\mathcal{P}\!\mathit{olarized}}
\def\Complexesstack{\mathcal{C}\!\mathit{omplexes}}
% \Pic is the operator that assigns to X its picard group, usage \Pic(X)
% \Picardstack_{X/B} denotes the Picard stack of X over B
% \Picardfunctor_{X/B} denotes the Picard functor of X over B
\def\Pic{\mathop{\mathrm{Pic}}\nolimits}
\def\Picardstack{\mathcal{P}\!\mathit{ic}}
\def\Picardfunctor{\mathrm{Pic}}
\def\Deformationcategory{\mathcal{D}\!\mathit{ef}}


\begin{document}

\title{Yoga}
\maketitle

\phantomsection
\label{section-phantom}

\tableofcontents

\section{Scope and sources}
\label{section-scope-and-sources}

This chapter has two parts. The first three
sections are short notes in English on the
classical background of yoga: the eight limbs,
the Yoga Sūtras of Patañjali, and the place of
yoga in the Mahābhārata. They follow
\cite{bryant-yoga-sutras} and
\cite{van-buitenen-mahabharata-1}.

The last section is a manual of Haṭha Yoga
postures in Spanish, written in 2023,
following B. K. S. Iyengar's Light on Yoga
\cite{iyengar-light-on-yoga}. It describes
about 80 postures and their variants. For
each posture it gives the instructions to
enter and leave it, the actions to perform
once in it, the benefits attributed to it
and, in some cases, the conditions under
which its practice is not recommended. The
statements attributed to Iyengar in the text
refer to that book.

The manual is kept in Spanish. It was lightly
edited for publication: the photographs of
the original are omitted, some of their
captions were turned into short lists of
variants, and sentences that only pointed to
a picture were removed. The groups of
postures of the original are the subsections
of the last section; within a subsection,
each posture has its own heading.

The manual describes a practice. The benefits
and contraindications listed are those given
in the manual's sources; they are not medical
claims and do not replace medical advice.

\section{The eight limbs of yoga}
\label{section-eight-limbs}

The Yoga Sūtras of Patañjali (see Section
\ref{section-yoga-sutras}) describe the
practice of yoga as having eight limbs, in
Sanskrit aṣṭāṅga (aṣṭa = eight, aṅga = limb).
They are listed in sūtra II.29
\cite{bryant-yoga-sutras}:
\begin{enumerate}
\item Yama: abstentions, rules of conduct
towards others.
\item Niyama: observances, rules of personal
discipline.
\item Āsana: posture.
\item Prāṇāyāma: regulation of the breath.
\item Pratyāhāra: withdrawal of the senses
from their objects.
\item Dhāraṇā: concentration, fixing the mind
on one place or object.
\item Dhyāna: meditation, the uninterrupted
continuation of that attention.
\item Samādhi: meditative absorption.
\end{enumerate}

The five yamas are listed in sūtra II.30:
\begin{enumerate}
\item Ahiṃsā: nonviolence.
\item Satya: truthfulness.
\item Asteya: not stealing.
\item Brahmacarya: celibacy.
\item Aparigraha: absence of greed, not
accumulating possessions.
\end{enumerate}

The five niyamas are listed in sūtra II.32:
\begin{enumerate}
\item Śauca: cleanliness.
\item Santoṣa: contentment.
\item Tapas: austerity.
\item Svādhyāya: study, in particular the
study of sacred texts.
\item Īśvarapraṇidhāna: devotion to Īśvara,
the Lord.
\end{enumerate}

The first five limbs are treated in the
second chapter of the Yoga Sūtras and the
last three at the beginning of the third.
The last three taken together are called
saṃyama (sūtra III.4).

Posture occupies little space in the Yoga
Sūtras: three sūtras, II.46 to II.48. Sūtra
II.46 says that posture should be steady and
comfortable (sthira-sukham āsanam). The
detailed systems of postures belong to the
later tradition of Haṭha Yoga. The manual in
Section \ref{section-posturas} concerns only
this third limb.

\section{The Yoga Sūtras of Patañjali}
\label{section-yoga-sutras}

The Yoga Sūtras are the classical Sanskrit
text of the school of Yoga, one of the six
classical schools of Hindu philosophy. The
text is attributed to Patañjali. Its date is
uncertain: nothing reliable is known about
its author, and scholarly estimates of the
date vary by several centuries, from the last
centuries before the common era to the first
centuries of it \cite{bryant-yoga-sutras}.

A sūtra is a very short aphorism, composed to
be memorized and to be explained by a teacher
or a commentary. The text has 195 or 196
sūtras, depending on the recension, in four
chapters (pāda):
\begin{enumerate}
\item Samādhi Pāda, on meditative absorption.
\item Sādhana Pāda, on practice. It contains
the eight limbs of Section
\ref{section-eight-limbs}.
\item Vibhūti Pāda, on the powers that are
said to arise from the practice.
\item Kaivalya Pāda, on liberation.
\end{enumerate}

The second sūtra of the first chapter, I.2,
defines yoga:
\begin{center}
yogaś citta-vṛtti-nirodhaḥ.
\end{center}
Here citta is the mind, vṛtti means
fluctuation or changing state, and nirodha
means cessation or stilling. The manual below
quotes this sūtra in Spanish: ``Yoga es el
cese de las fluctuaciones de la mente''.

Because the sūtras are so brief, the text has
always been read with a commentary. The
oldest and most influential one is the
Bhāṣya attributed to Vyāsa; the later
commentators comment on the sūtras together
with it \cite{bryant-yoga-sutras}.

\section{Yoga in the Mahābhārata}
\label{section-yoga-in-the-mahabharata}

The Mahābhārata is one of the two great
Sanskrit epics of India; the other one is the
Rāmāyaṇa. Tradition attributes it to the sage
Vyāsa. It is divided into eighteen books
(parvan), and its main story is the conflict
between two groups of cousins, the Kauravas
and the Pāṇḍavas, over the succession to the
throne, which ends in a great war
\cite{van-buitenen-mahabharata-1}.

The Bhagavad Gītā is part of the sixth book
of the epic, the Bhīṣma Parva. It is a
dialogue in eighteen chapters between the
Pāṇḍava prince Arjuna and Kṛṣṇa, his
charioteer, placed just before the battle
begins. Arjuna does not want to fight against
his own relatives and teachers, and Kṛṣṇa
answers him.

In the Gītā the word yoga is used for several
paths, among them:
\begin{enumerate}
\item Karma yoga, the path of action
performed without attachment to its results.
\item Jñāna yoga, the path of knowledge.
\item Bhakti yoga, the path of devotion.
\end{enumerate}

The manual below refers to the epic once. The
posture Bharadvājāsana is named after the
sage Bharadvāja, the father of Droṇa. Droṇa
is the teacher who instructs both the
Kauravas and the Pāṇḍavas in the use of
weapons; in the war he fights on the side of
the Kauravas.

\section{Posturas}
\label{section-posturas}

Este manual describe una práctica. Los
beneficios y contraindicaciones que se
enumeran son los que dan las fuentes del
manual; no son afirmaciones médicas y no
sustituyen el consejo médico.

\subsection{Acerca de este manual}
\label{section-acerca-de-este-manual}

Este manual es una guía de Haṭha Yoga con
alrededor de 80 posturas con sus variantes.
Incluimos para cada una las instrucciones
para entrar y salir de la postura, las
acciones a realizar una vez en ella, los
beneficios que se le atribuyen y, en algunos
casos, en qué condiciones no se recomienda
practicarla.

Sigue el libro Light on Yoga de B. K. S.
Iyengar \cite{iyengar-light-on-yoga}.

Muchas veces, las instrucciones para entrar a
las posturas están escritas pensando en la
postura completa. Esto puede parecer un
desafío enorme, sobre todo para quien
comienza a practicar. Una recomendación
importantísima es \textbf{no ir más allá de
nuestro límite}. Nuestro objetivo simplemente
es comprender cómo funcionan las posturas:
cómo entrar a ellas, qué partes del cuerpo
están involucradas, cómo profundizar, etc.
Más que concentrarnos en lograr las posturas
en su máxima expresión, debemos
concentrarnos: ``Yoga es el cese de las
fluctuaciones de la mente''. Ése es nuestro
objetivo.

Mencionamos aquí algunas posturas que no
están incluidas en este manual:
\begin{itemize}
\item Balāsana.
\item Gato-Vaca (Bidālāsana o Mārjārīāsana).
\item Paloma (Eka Pāda Rājakapotāsana).
\item Hanumanāsana.
\item Parivṛtta Jānu Sīrsāsāna.
\item Sīrsāsana II y su transición a Bakāsana
(y demás variantes).
\item Ardha Baddha Padmottānāsana y su
transición a Vātāyanāsana.
\item Agni Stambhāsana.
\item Vīrabhadrāsana III.
\item Y las que aparecen en los Saludos al
Sol: Plancha (Phalakāsana), Ānjaneyāsana y su
variante Lizard Pose (Utthana Pristhāsana),
Ashtangāsana (ocho puntos) y Ardha
Uttānāsana.
\end{itemize}

Dwi Pāda Viparīta Daṇḍāsana y Setu Bhanda
Sarvāngāsana están clasificadas tanto como
posturas restaurativas como posturas de
extensión: como presentamos versiones con
props, nos quedamos con la primera opción en
ambos casos.

\subsection{Posturas preparatorias básicas}
\label{section-posturas-preparatorias}

\subsubsection*{Supta Pādānguṣṭasāsana I, II, III, IV}

Supta = acostado. Pādā = pie. Anguṣṭa = dedo
gordo del pie.

\emph{Instrucciones.}
\begin{enumerate}
\item Comenzar acostado, con ambas piernas
estiradas y activas.
\item Inhalar y levantar pierna izquierda en
90º sobre el piso. Pierna derecha estirada,
pie en punta y mano derecha descansando sobre
el muslo derecho.
\item Con dedos índice y medio de la mano
izquierda tomar el dedo gordo del pie
izquierdo. Respirar 3 o 4 veces.
\item En una exhalación levantar la frente
hacia la rodilla izquierda. Jalar la pierna
con el brazo izquierdo sin doblar la rodilla.
20 segundos sin doblar rodilla.
\item En una inhalación devolver cabeza al
piso y pierna izquierda en 90º sin soltar el
agarre de la mano izquierda. Ésta fue Supta
Pādānguṣṭasāsana IV.
\item En una exhalación doblar rodilla y
llevar el pie izquierdo hacia el hombro
derecho. Doblar el codo izquierdo y levantar
la cabeza para llevar el brazo izquierdo por
detrás de la cabeza. Barbilla al pecho y
respirar profundo.
\item En una inhalación devolver cabeza al
piso y el brazo izquierdo por delante de la
cabeza. Luego llevar cabeza al piso y pierna,
sin soltar el agarre, de regreso a la
perpendicular. Todo este tiempo la pierna
derecha bien estirada y la mano derecha sobre
el muslo derecho. Hasta ahora fue Supta
Pādānguṣṭasāsana II.
\item En una exhalación llevar el pie
izquierdo al piso rotando externamente la
cadera (hacia el lado izquierdo). Llevar el
brazo a la altura de los hombros jalando pie
izquierdo hacia hombro. 20 segundos
respirando normalmente.
\item En una inhalación devolver la pierna a
los 90º sin soltar el dedo gordo y sin doblar
la rodilla, y sin alterar la posición
estirada de la pierna derecha y de la mano
derecha sobre el muslo derecho. Éste fue
Supta Pādānguṣṭasāsana II. Iyengar termina
aquí la secuencia de Supta Pādānguṣṭasāsana.
\item En una exhalación llevar el pie
izquierdo al piso rotando internamente la
cadera (hacia el lado derecho). Llevar el
brazo a la altura de los hombros jalando pie
izquierdo hacia hombro. 20 segundos
respirando normalmente.
\item En una inhalación devolver la pierna a
los 90º sin soltar el dedo gordo y sin doblar
la rodilla, y sin alterar la posición
estirada de la pierna derecha y de la mano
derecha sobre el muslo derecho. Éste fue
Supta Pādānguṣṭasāsana III.
\item En una exhalación soltar el pie
izquierdo y devolver la pierna izquierda al
piso junto con la derecha.
\item Respirar profundamente y hacer el otro
lado.
\end{enumerate}

\emph{Acciones.}
\begin{itemize}
\item Si es difícil tomar el dedo gordo del
pie sin doblar rodilla y sin flexionar la
espalda usar cinturón en la mitad del pie. En
Supta Pādānguṣṭasāsana I se toma con ambas
manos, una en cada lado del cinturón. En
Supta Pādānguṣṭasāsana II y III el cinturón
se toma con una sola mano, la que está en la
dirección en la que baja la pierna.
\item Si es difícil mantener la pierna de
abajo en el piso se puede pedir ayuda a un
amigo que haga presión hacia abajo empujando
el muslo. También se puede hacer presión a
una pared.
\end{itemize}

\emph{Beneficios.} Ayuda a activar piernas.
La sangre fluye en las piernas y los nervios
de la cadera se rejuvencen. Alivia problemas
en ciática.

\subsubsection*{Supta Tādāsana con Ūrdhva Hastāsana}

Supta: acostado. Tāda = montaña. Ūrdhva =
arriba, alto. Hasta = mano.

\emph{Instrucciones.}
\begin{enumerate}
\item Acostarse en el piso con las piernas
extendidas y juntas. Activar rótulas, muslos
y flexionar pies.
\item Activar abdomen y meter coxis.
\item Llevar lentamente los brazos hacia
atrás bien estirados.
\end{enumerate}

\emph{Acciones.} Una variante de la posición
de los brazos es tomándose de los codos:
levantar los bazos perpendiculares al piso y
tomarse de los codos con las manos. Llevar
los brazos entrelazados de esta forma por
detrás de la cabeza. Después de unas
respiraciones regresar los brazos
entrelazados frente al pecho y cambiar el
entrelazado. Llevar nuevamente brazos hacia
atrás y respirar.

\emph{Beneficios.} La columna se estabiliza y
se extiende. Descansan los músculos en la
espalda. La mente se relaja fácilmente. Puede
considerarse una postura preparativa para
Sūryanamaskar.

\subsubsection*{Pávana Muktāsana}

Pavana = viento. Mukta = liberación.

\emph{Instrucciones.}
\begin{enumerate}
\item Comenzar en Savāsana.
\item Flexionar rodilla derecha hasta
entrelazar dedos en la espinilla.
\item Acercar rodilla derecha a hombro
derecha. Respirar.
\item Repetir con rodilla izquierda.
\item Repetir con ambas rodillas. En este
caso tratar de tomar codos con ambas manos
mientras la espalda está completamente
apoyada en el piso.
\item Estirar ambas piernas y volver a
Savāsana.
\end{enumerate}

\emph{Acciones.}
\begin{itemize}
\item En la exhalación acercar más la rodilla
hacia el hombro. En la postura completa hay
contacto.
\item Se puede explorar con movimientos
pequeños y conscientes qué partes del
cuadriceps y del abdomen se activan.
\item Variante: levantar la frente del piso y
acercar a rodilla.
\item Si es complicado entrelazar manos se
puede:
\begin{itemize}
\item Rotar la rodilla externamente.
\item Usar cinturón.
\end{itemize}
\end{itemize}

\emph{Beneficios.} Liberación de gases en el
estómago. Alivia procesos digestivos.

\emph{Contraindicaciones.} Cirugías
abdominales, hernias. Problemas en columna,
ciática. Embarazo.

\subsection{Posturas de pie}
\label{section-posturas-de-pie}

\medskip\noindent\emph{Grupo 1: Posturas con
eje completamente vertical}

\subsubsection*{Tādāsana}

(También llamada Samasthiti) Tāda = montaña.
Sama = erguido, recto, inmóvil. Sthiti =
inmóvil.

\emph{Instrucciones.}
\begin{enumerate}
\item Pararse de pie con los pies juntos, los
talones y los dedos gordos de los pies se
tocan. Abrir los tres arcos del pie y separar
los dedos.
\item Activar rótulas, muslos y pantorrillas.
Coxis adentro.
\item Estómago adentro, espalda estirada,
tracción de columna, proyectar esternón,
barbilla ligeramente al pecho, coronilla al
cielo.
\item Los brazos a los lados del cuerpo.
Palmas mirando hacia los muslos.
\end{enumerate}

\emph{Beneficios.} Estar conscientes de cómo
es nuestra postura cuando estamos de pie es
fundamental para la vida. Esta postura nos
ayuda a corregir el modo en que está nuestro
cuerpo más allá de la práctica de yoga, y
esto puede generar cambios importantes en
nosotros. Una mala postura de pie puede
generar grandes desequilibrios y deformidades
en el cuerpo a largo plazo.

\emph{Variantes.} Variantes de la posición de
los brazos: Ūrdhva Hastāsana, Ūrdhva
Baddhanguliyāsana, Namaskarāsana, Paschima
Namaskarāsana y Gomukhāsana.

\subsubsection*{Utkatāsana}

Utkata = poderoso, vigoroso.

\emph{Instrucciones.}
\begin{enumerate}
\item Comenzar en Tadāsana. Levantar los
brazos por arriba de la cabeza y juntar
palmas.
\item Exhalar y doblar rodillas hasta que los
muslos estén paralelos al piso manteniendo
columna bien estirada.
\item Respirar profundo. 30 segundos es más
que suficiente para esta postura.
\item Inhalar y estirar rodillas con los
brazos aun estirados. Regresar brazos abajo
para volver a Tadāsana.
\end{enumerate}

\emph{Beneficios.} Ayuda a relajar
contracturas en hombros y corrige
deformidades menores en las piernas. Los
tobillos se fortalecen y los muslos se
desarrollan equilibradamente. El diafragma se
levanta y esto masajea el corazón. Tonifica
abdomen y abre pecho.

\subsubsection*{Garuḍāsana}

Garuḍa = águila. Se representa como el
vehículo de Viṣṇu con cara blanca, alas rojas
y cuerpo de oro.

\emph{Instrucciones.}
\begin{enumerate}
\item Comenzar en Tadāsana.
\item Flexionar la rodilla derecha y cruzar
la pierna izquierda sobre la derecha. Si
posible, entrelazar el pie izquierdo en el
tobillo derecho.
\item Cruzar los brazos, izquierdo por
arriba, de tal forma que los codos se toquen
y los antebrazos estén perpendiculares al
piso. Las palmas se tocan.
\item Respirar profundo por 15 o 20 segundos.
\item Deshacer el entrelazado de brazos,
luego de piernas y volver a Tadāsana.
\item Hacer el otro lado.
\end{enumerate}

\emph{Beneficios.} Ayuda a fortalecer
tobillos y relaja hombros. Previene calambres
en las pantorrillas.

\medskip\noindent\emph{Grupo 2. Posturas con
piernas separadas de manera horizontal}

\subsubsection*{Uttihita Hasta Pādāsana}

Uttihita = extendido, estirado. Hasta = mano.
Pādā = pierna

\emph{Instrucciones.}
\begin{enumerate}
\item Comenzar en Tadāsana.
\item Abrir el compás de las piernas con un
salto o dando un paso. Los bordes externos de
los pies están paralelos a los bordes del
tapete.
\item Levantar los brazos paralelos al piso,
palmas mirando hacia el piso.
\item Realizar las acciones y respirar
profundo.
\item Para salir, exhalar y doblar rodillas.
Se puede salir de un brinco o dar un paso
para juntar piernas y volver a Tadāsana.
\end{enumerate}

\emph{Acciones.}
\begin{itemize}
\item Estirar columna, meter coxis .
\item Rotar la articulación del fémur hacia
adentro.
\item Juntar omóplatos.
\item Dependiendo de la flexibilidad de las
caderas de la persona se puede abrir menos o
más el compás de las piernas.
\end{itemize}

\emph{Beneficios.} Sirve para ejercitar los
músculos internos de los muslos. Corrige
curvatura lumbar, activa espalda.

\emph{Contraindicaciones.}
\begin{itemize}
\item En caso de dolencias severas en cadera
mejor no trabajar esta postura.
\item En caso de dolencias de hombros y
cuello poner manos en caderas.
\item En caso de embarazo, menstruación,
dolor en rodillas o columna no entrar
saltando.
\end{itemize}

\subsubsection*{Pārśva Hasta Pādāsana}

Pārśva = lado o flanco. Hasta = mano. Pādā =
pierna.

\emph{Instrucciones.}
\begin{enumerate}
\item Comenzar en Tadāsana.
\item Abrir el compás de las piernas con un
salto o dando un paso. Los bordes externos de
los pies están paralelos a los bordes del
tapete.
\item Levantar los brazos paralelos al piso,
palmas mirando hacia el piso.
\item Girar el pie izquierdo hacia el centro
del cuerpo de tal forma que el talón quede
hacia afuera en un ángulo de 30º. Activar
rótula y toda la pierna izquierda. Girar la
pierna derecha desde la articulación abriendo
la cadera hasta que el pie quede en ángulo de
90º.
\item Respirar y realizar las acciones.
\item Regresar a Uttihita Hasta Pādāsana como
se entró: primero regresar pierna derecha y
luego pierna izquierda.
\item Hacer el otro lado.
\item Salir de Uttihita Hasta Pādāsana por
Tadāsana
\end{enumerate}

\emph{Acciones.}
\begin{itemize}
\item Rotar muslo izquierdo hacia afuera.
\item Comprobar que el pie izquierdo está en
línea con la mitad arco interno del pie
derecho..
\item Estirar y activar bien ambas piernas.
\item Comprobar que el peso esté bien
distribuido en ambas piernas.
\item Meter el coxis.
\item Pubis, ombligo, esternón y nariz en una
línea vertical perpendicular la piso, mirada
hacia en frente.
\item Elevar el torso y proyectar esternón.
\item Comprobar que los brazos están bien
estirados paralelos al piso y las palmas
mirando al piso.
\item Estirar músculos del cuello, relajar el
rostro y respirar profundo.
\item Dependiendo de la flexibilidad de las
caderas de la persona se puede abrir menos o
más el compás de las piernas.
\end{itemize}

\emph{Beneficios.} Se nivelan las caderas.
Abre el pecho. Fortalece rodillas, piernas y
tobillos. Mejora estabilidad y concentración.

\emph{Contraindicaciones.}
\begin{itemize}
\item En caso de dolencias severas en cadera
mejor no trabajar esta postura.
\item En caso de dolencias de hombros y
cuello poner manos en caderas.
\item En caso de embarazo, menstruación,
dolor en rodillas o columna no entrar
saltando.
\end{itemize}

\subsubsection*{Uttihita Trikoṇāsana}

Uttihita = extendido, estirado. Trikoṇa =
triángulo (Tri = tres, kona = ángulo). Es la
postura del triángulo extendido.

\emph{Instrucciones.}
\begin{enumerate}
\item Comenzar en Tadāsana.
\item Separar las piernas con un salto o
dando un paso. Dependiendo de la flexibilidad
en caderas, la distancia es alrededor de un
metro. Levantar los brazos estirados
paralelos al piso con las palmas mirando
hacia abajo.
\item Girar el pie izquierdo hacia el centro
del cuerpo de tal forma que el talón quede
hacia afuera en un ángulo de 30º. Activar
rótula y toda la pierna izquierda. Girar la
pierna derecha desde la articulación abriendo
la cadera hasta que el pie quede en ángulo de
90º.
\item Girar ligeramente la
cadera izquierda hacia el pie derecho.
\item En una exhalación acercar el torso al
piso y luego la mano. En la postura completa
la palma toca el piso por fuera del tobillo
derecho.
\item Estirar el brazo izquierdo hacia el
cielo el línea con el brazo derecho. Mirada
hacia la mano izquierda.
\item Mantener las piernas bien activas y
respirar profundo entre 30 segundos y un
minuto.
\end{enumerate}

\emph{Acciones.}
\begin{itemize}
\item Comprobar que ambos brazos estén en una
línea.
\item Juntar omóplatos para generar una
ligera curvatura en espalda alta.
\item Comprobar que todo el cuerpo esté en un
plano perpendicular al piso.
\item Mantener pubis, ombligo, esternón y
nariz en una misma línea.
\item Cóxis adentro.
\item Empujar glúteo derecho hacia adelante.
\item Rostro relajado y respiración profunda.
\end{itemize}

\emph{Salida.}
\begin{enumerate}
\item Inhalar y levantar la mano derecha
junto con el torso para terminar en la
posición del paso 5.
\item Para hacer el otro lado se pueden girar
la cadera y ambos pies para entrar en la
postura.
\item Saltar o dar un paso de regreso a
Tadāsana.
\end{enumerate}

\emph{Variantes/Adaptaciones.}
\begin{itemize}
\item Microflexionar rodilla de adelante.
\item Mirada al suelo o hacia en frente.
\item Mano de arriba en la cadera.
\item Usar bloque por fuera del pie en
cualquiera de las tres alturas.
\end{itemize}

\emph{Beneficios.}
\begin{itemize}
\item Estira y fortalece muslos, rodillas y
pantorrillas.
\item Estimula abromen y sistema digestivo.
\item Alivia dolores de espalda relacionados
al embarazo.
\end{itemize}

\emph{Contraindicaciones.} Diarrea, dolor de
cabeza, presión arterial baja.

\subsubsection*{Vīrabhadrāsana II}

Vīrabhadra = un héroe que surgió de un
cabello de Śiva.

\emph{Instrucciones.}
\begin{enumerate}
\item Comenzar en Tadāsana.
\item Abrir el compás de las piernas con un
salto o dando un paso. Los bordes externos de
los pies están paralelos a los bordes del
tapete.
\item Levantar los brazos paralelos al piso,
palmas mirando hacia el piso.
\item Girar el pie izquierdo hacia el centro
del cuerpo de tal forma que el talón quede
hacia afuera en un ángulo de 30º. Activar
rótula y toda la pierna izquierda. Girar la
pierna derecha desde la articulación abriendo
la cadera hasta que el pie quede en ángulo de
90º.
\item Doblar la pierna derecha hasta llegar a
un ángulo recto, con el muslo paralelo al
piso alargando la musculatura de la ingle.
\item Mirada sobre el brazo derecho.
\item Respirar profundo y sostener la
postura.
\item Para salir simplemente estirar la
pierna derecha para volver a Pārśva Hasta
Pādāsana y luego a Tadāsana.
\item Hacer el otro lado.
\end{enumerate}

\emph{Acciones.}
\begin{itemize}
\item Cóxis adentro.
\item Comprobar que la rodilla de adelante no
vaya más allá del tobillo: deben estar en el
mismo plano.
\item Estirar columna, proyectar esternón.
\item Estira la parte posterior del cuello.
\item Estirar los brazos.
\end{itemize}

\emph{Variantes/Adaptaciones.}
\begin{itemize}
\item Durante el embarazo, manos en cadera.
\item Si hay falta de flexibilidad en
caderas, disminuir distancia entre pies.
\item Si hay tensión en cuello, mirada al
frente.
\end{itemize}

\emph{Beneficios.}
\begin{itemize}
\item Fortalece y da flexibilidad en piernas
y tobillos.
\item Perfeccionar esta postura ayuda como
preparación para las flexiones hacia
adelante.
\item Alivia dolor de espalda.
\item Es una buena postura para practicar la
concentración durante la práctica.
\end{itemize}

\emph{Contraindicaciones.} Diarrea, dolor de
cabeza, presión arterial baja.

\subsubsection*{Uttihita Pārśvkoṇāsana}

Uttihita = extendido, estirado. Pārśva = lado
o flanco. Kona = ángulo.

\emph{Instrucciones.}
\begin{enumerate}
\item Comenzar en Tadāsana.
\item Abrir el compás de las piernas con un
salto o dando un paso. Los bordes externos de
los pies están paralelos a los bordes del
tapete.
\item Levantar los brazos paralelos al piso,
palmas mirando hacia el piso.
\item Girar el pie izquierdo hacia el centro
del cuerpo de tal forma que el talón quede
hacia afuera en un ángulo de 30º. Activar
rótula y toda la pierna izquierda. Girar la
pierna derecha desde la articulación abriendo
la cadera hasta que el pie quede en ángulo de
90º.
\item Doblar la pierna derecha hasta llegar a
un ángulo recto, con el muslo paralelo al
piso alargando la musculatura de la ingle. La
pierna izquierda queda bien estirada y el pie
izquierdo en el piso.
\item Llevar la palma derecha al piso por
afuera del pie. La axila derecha queda sobre
la rodilla.
\item Estirar el brazo izquierdo por arriba
de la oreja izquierda. Mirada hacia el techo
por dentro de la axila izquierda.
\end{enumerate}

\emph{Acciones.}
\begin{itemize}
\item Levantar el pecho y arquear ligeramente
espalda alta. El torso y la cadera deben
estar en el mismo plano.
\item Estirar la columna sintiendo separación
en las vértebras (tracción de columna) y en
las costillas.
\item Estirar por completo el costado del
brazo que está arriba.
\item Es sumamente útil usar un bloque en la
mano de abajo para mantener la columna recta
y concentrarse en el estiramiento.
\end{itemize}

\emph{Salida.}
\begin{enumerate}
\item Inhalando estirar la pierna derecha y
levantar los brazos hasta la horizontal.
\item Rotar los pies para hacer el otro lado.
\item Para salir dar un paso o un brinco para
terminar en Tadāsana.
\end{enumerate}

\emph{Variantes/Adaptaciones.}
\begin{itemize}
\item Microflexionar rodilla de adelante.
\item Mirada al suelo o hacia en frente.
\item Mano de arriba en la cadera.
\item Usar bloque por fuera del pie en
cualquiera de las tres alturas.
\end{itemize}

\emph{Beneficios.}
\begin{itemize}
\item Tonifica tobillos, rodillas y muslos.
\item Alivia dolores de ciática y artritis.
\item Ayuda a los procesos de eliminación en
el recto.
\end{itemize}

\medskip\noindent\emph{Grupo 3. Posturas con
piernas separadas una detrás de la otra}

\subsubsection*{Vimanasana}

Viman = vehículo volador.

\emph{Instrucciones.}
\begin{enumerate}
\item Comenzar en Tadāsana.
\item Abrir el compás de las piernas con un
salto o dando un paso. Los bordes externos de
los pies están paralelos a los bordes del
tapete.
\item Levantar los brazos paralelos al piso,
palmas mirando hacia el piso.
\item Girar el pie izquierdo hacia el centro
del cuerpo de tal forma que el talón quede
hacia afuera en un ángulo de 30º. Activar
rótula y toda la pierna izquierda. Girar la
pierna derecha desde la articulación abriendo
la cadera hasta que el pie quede en ángulo de
90º.
\item Llevar las manos a las caderas y girar
el torso hasta mirar hacia la pierna derecha.
\item Doblar la pierna derecha hasta llegar a
un ángulo recto, con el muslo paralelo al
piso alargando la musculatura de la ingle.
\item Estirar los brazos a los lados con las
palmas mirando hacia abajo.
\end{enumerate}

\emph{Acciones.}
\begin{itemize}
\item Cuidar que la rodilla de la pierna de
adelante esté justo arriba del pie y el muslo
en la misma línea.
\item Pierna de atrás bien estirada y rótula
activa.
\item Hacer conciente el arco externo del pie
de adelante.
\item Coxis adentro.
\item Elevar el torso desde el pubis, estirar
espalda y proyectar esternón.
\item Omóplatos se juntan, abrir el pecho.
\item Externder la parte posterior del cuello
\item Relajar el rostro y respirar
profundamente.
\end{itemize}

\emph{Salida.}
\begin{enumerate}
\item Con una exhalación bajar los brazos de
regreso a la cadera.
\item Con una inhalación estirar la pierna
derecha y devolver el talón izquierdo al
piso.
\item Regresar el torso al centro y los pies
paralelos.
\item Hacer el otro lado.
\item Para salir, con un brinco o un paso
volver a Tadāsana.
\end{enumerate}

\emph{Variantes/Adaptaciones.}
\begin{itemize}
\item Una variante es con las palmas mirando
hacia arriba.
\item Durante el embarazo o si hay tensión en
hombros: manos en cadera.
\end{itemize}

\emph{Beneficios.} Fortalece las piernas,
alarga espalda, libera la tensión en hombros
y cadera. Estabilidad, tierra, calma en la
mente.

\subsubsection*{Vīrabhadrāsana I}

Vīrabhadra = un héroe que surgió de un
cabello de Śiva.

\emph{Instrucciones.}
\begin{enumerate}
\item Comenzar en Tadāsana.
\item Abrir el compás de las piernas con un
salto o dando un paso. Los bordes externos de
los pies están paralelos a los bordes del
tapete.
\item Levantar los brazos paralelos al piso,
palmas mirando hacia el piso.
\item Girar el pie izquierdo hacia el centro
del cuerpo de tal forma que el talón quede
hacia afuera en un ángulo de 30º. Activar
rótula y toda la pierna izquierda. Girar la
pierna derecha desde la articulación abriendo
la cadera hasta que el pie quede en ángulo de
90º.
\item Llevar las manos a las caderas y girar
el torso hasta mirar hacia la pierna derecha.
\item Elevar las manos por arriba de la
cabeza y juntar palmas.
\item Doblar la pierna derecha hasta llegar a
un ángulo recto, con el muslo paralelo al
piso alargando la musculatura de la ingle.
\end{enumerate}

\emph{Acciones.}
\begin{itemize}
\item Rodilla de la pierna derecha en línea
con la pierna izquierda.
\item Pierna izquierda bien estirada con
rótula activa.
\item Comprobar que el peso esté bien
distribuido en ambas piernas.
\item Coxis adentro.
\item Juntar omóplatos, proyectar esternón y
estirar espalda.
\item Para profundizar, mirada hacia el cielo
y arquear espalda alta.
\end{itemize}

\emph{Salida.}
\begin{enumerate}
\item Inhalar y estirar la pierna derecha.
\item Manos a cadera y devolver torso y pies
al centro.
\item Hacer el otro lado.
\item Para salir, con un brinco o un paso
volver a Tadāsana.
\end{enumerate}

\emph{Beneficios.}
\begin{itemize}
\item Estira y fortalece la columna
vertebral.
\item Fortalece rodillas.
\item Aumenta flexibilidad de cadera.
\item Mejora sistema circulatorio.
\item Mejora la capacidad pulmonar.
\item Fortaleza, determinación y constancia.
\end{itemize}

\emph{Adaptaciones.} Si hay tensión hombros,
cuello, caderas, o si hay embarazo, llevar
las manos a la cadera.

\medskip\noindent\emph{Grupo 4. Posturas con
torsión}

\subsubsection*{Parivṛtta Trikoṇāsana}

Parivṛtta = girado, torcido. Trikoṇa =
triángulo (Tri = tres, kona = ángulo). Es la
postura del triángulo torcido.

\emph{Instrucciones.}
\begin{enumerate}
\item Comenzar en Tadāsana.
\item Separar las piernas con un salto o
dando un paso. Dependiendo de la flexibilidad
en caderas, la distancia es alrededor de un
metro. Levantar los brazos estirados
paralelos al piso con las palmas mirando
hacia abajo.
\item Girar el pie izquierdo hacia el centro
del cuerpo de tal forma que el talón quede
hacia afuera en un ángulo de 30º. Activar
rótula y toda la pierna izquierda. Girar la
pierna derecha desde la articulación abriendo
la cadera hasta que el pie quede en ángulo de
90º.
\item Girar ligeramente la
cadera izquierda hacia el pie derecho.
\item En una exhalación acercar el torso al
piso y la mano izquierda va por afuera del
pie derecho.
\item Estirar el brazo derecho hacia el techo
en línea con el brazo izquierdo. Mirada hacia
la mano derecha.
\end{enumerate}

\emph{Acciones.}
\begin{itemize}
\item Piernas bien estiradas con rótulas
activas y muslos activos.
\item Columna vertebral bien estirada
proyectando esternón.
\item Ligero arco en espalda alta.
\item Coxis adentro.
\item Rostro relajado y respiración profunda.
\item Conectar con la torsión de columna.
\end{itemize}

\emph{Salida.}
\begin{enumerate}
\item En una inhalación levantar la mano
izquierda del piso y devolver el tronco a su
porición normal.
\item Respirar profundo y rotar los pies a la
posición correspondiente para hacer el otro
lado.
\item Para salir, con un brinco o un paso
volver a Tadāsana.
\end{enumerate}

\emph{Variantes y contraindicaciones.}
\begin{itemize}
\item Usar un bloque ayuda mucho a mantener
la columna estirada y recta.
\item Para hacer la postura menos intensa,
descansar la mano arriba sobre el coxis y
llevar la mirada al suelo.
\item No realizar durante el embarazo.
\end{itemize}

\emph{Beneficios.}
\begin{itemize}
\item Regula sistema digestivo y otros
órganos internos.
\item Fortalece piernas.
\item Mejora equilibrio.
\item Fortalece músculos abdominales.
\item Alivia dolores de espalda
\end{itemize}

\subsubsection*{Parivṛtta Pārśvkoṇāsana}

Parivṛtta = girado, torcido. Pārśva = lado o
flanco. Kona = ángulo.

\emph{Instrucciones.}
\begin{enumerate}
\item Comenzar en Tadāsana.
\item Abrir el compás de las piernas con un
salto o dando un paso. Los bordes externos de
los pies están paralelos a los bordes del
tapete.
\item Levantar los brazos paralelos al piso,
palmas mirando hacia el piso.
\item Girar el pie izquierdo hacia el centro
del cuerpo de tal forma que el talón quede
hacia afuera en un ángulo de 30º. Activar
rótula y toda la pierna izquierda. Girar la
pierna derecha desde la articulación abriendo
la cadera hasta que el pie quede en ángulo de
90º.
\item En una exhalación rotar el torso para
llevar la palma izquierda al piso por fuera
del pie derecho.
\item Profundizar en la torsión y estirar el
brazo derecho por arriba de la oreja derecha.
Mirada hacia el techo por dentro de la axila
derecha.
\end{enumerate}

\emph{Acciones.}
\begin{itemize}
\item Levantar el pecho y arquear ligeramente
espalda alta.
\item Estirar la columna sintiendo separación
en las vértebras (tracción de columna) y en
las costillas.
\item Es sumamente útil usar un bloque en la
mano de abajo para mantener la columna recta
y concentrarse en el estiramiento.
\item Cadera izquierda adelante.
\item Piernas bien estiradas y rótulas
activas.
\end{itemize}

\emph{Salida.}
\begin{enumerate}
\item Inhalando salir de la torsión y estirar
la pierna derecha, levantar los brazos hasta
la horizontal.
\item Rotar los pies para hacer el otro lado.
\item Para salir dar un paso o un brinco para
terminar en Tadāsana.
\end{enumerate}

\emph{Variantes y contraindicaciones.}
\begin{itemize}
\item Usar un bloque ayuda mucho a mantener
la columna estirada y recta.
\item Para hacer la postura menos intensa,
descansar la mano arriba sobre el coxis y
llevar la mirada al suelo.
\item No realizar durante el embarazo.
\end{itemize}

\emph{Beneficios.}
\begin{itemize}
\item Hay una gran estimulación en órganos
abdominales: facilita la digestión y
rejuvenece estos órganos.
\item Tonifica tobillos, rodillas y muslos.
\item Alivia dolores de ciática y artritis.
\item Mejora equilibrio.
\item Fortalece músculos abdominales.
\item Alivia dolores de espalda
\item Es más intensa que Parivṛtta
Trikoṇāsana y en general trae más beneficios
aunque en los isquitibiales no hay tanto
estiramiento.
\end{itemize}

\medskip\noindent\emph{Grupo 5. Posturas en
equilibrio}

\subsubsection*{Ardha Chandrāsana}

Ardha = medio. Chandra = luna.

\emph{Instrucciones.}
\begin{enumerate}
\item Comenzar en Tadāsana.
\item Separar las piernas con un salto o
dando un paso. Dependiendo de la flexibilidad
en caderas, la distancia es alrededor de un
metro. Levantar los brazos estirados
paralelos al piso con las palmas mirando
hacia abajo.
\item Girar el pie izquierdo hacia el centro
del cuerpo de tal forma que el talón quede
hacia afuera en un ángulo de 30º. Activar
rótula y toda la pierna izquierda. Girar la
pierna derecha desde la articulación abriendo
la cadera hasta que el pie quede en ángulo de
90º.
\item Girar ligeramente la
cadera izquierda hacia el pie derecho.
\item En una exhalación acercar el torso al
piso y la mano derecha va al piso por delante
del pie derecho. La pierna de atrás bien
estirada paralela al piso.
\item Brazo izquierdo en línea con el
derecho.
\item Mirada hacia la mano derecha.
\end{enumerate}

\emph{Acciones.}
\begin{itemize}
\item Meter cóccix.
\item Estirar bien las dos piernas.
\item Pies flexionados.
\item Espalda bien estirada, tracción de
columna y arco ligero en espalda alta.
\item Músculos del cuello bien estirados.
\end{itemize}

\emph{Salida.}
\begin{enumerate}
\item Regresar mirada al piso, brazo de
arriba al piso y pie de arriba al piso.
\item Devolver el torso a su posición normal
y girar los pies para hacer el otro lado.
\item Para salir dar un paso o un brinco para
terminar en Tadāsana.
\end{enumerate}

\emph{Variantes y contraindicaciones.}
\begin{itemize}
\item Usar un bloque ayuda mucho para
mantener una línea desde la pierna hasta la
coronilla, estirando bien la columna.
\item Llevar la mano de arriba al coxis y
mantener la mirada al piso ayudan a hacer más
ligera la postura.
\end{itemize}

\emph{Beneficios.} Ayuda a trabajar el
equilibrio, estirar espalda, desarrollar
fuerza en piernas. Mejora la concentración.

\subsubsection*{Parivṛtta Ardha Chandrāsana}

Parivṛtta = girado, torcido. Ardha = medio.
Chandra = luna.

\emph{Instrucciones.}
\begin{enumerate}
\item Comenzar en Tadāsana.
\item Separar las piernas con un salto o
dando un paso. Dependiendo de la flexibilidad
en caderas, la distancia es alrededor de un
metro. Levantar los brazos estirados
paralelos al piso con las palmas mirando
hacia abajo.
\item Girar el pie izquierdo hacia el centro
del cuerpo de tal forma que el talón quede
hacia afuera en un ángulo de 30º. Activar
rótula y toda la pierna izquierda. Girar la
pierna derecha desde la articulación abriendo
la cadera hasta que el pie quede en ángulo de
90º.
\item Girar ligeramente la
cadera izquierda hacia el pie derecho.
\item En una exhalación acercar el torso al
piso y la mano izquierda va al piso por
delante del pie derecho. La pierna de atrás
bien estirada paralela al piso.
\item Brazo derecho en línea con el izquierdo
hacia arriba.
\item Mirada al techo.
\end{enumerate}

\emph{Acciones.}
\begin{itemize}
\item Espalda bien estirada, columna en una
misma línea con la pierna izquierda.
\item Ligera curvatura en espalda alta.
\item Relajar el rostro y respirar profundo.
\end{itemize}

\emph{Salida.}
\begin{enumerate}
\item Regresar mirada al piso, brazo de
arriba al piso y pie de arriba al piso.
\item Devolver el torso a su posición normal
y girar los pies para hacer el otro lado.
\item Para salir dar un paso o un brinco para
terminar en Tadāsana.
\end{enumerate}

\emph{Variantes y contraindicaciones.}
\begin{itemize}
\item Usar un bloque ayuda mucho para
mantener una línea desde la pierna hasta la
coronilla, estirando bien la columna.
\item Llevar la mano de arriba al coxis y
mantener la mirada al piso ayudan a hacer más
ligera la postura.
\end{itemize}

\emph{Beneficios.} Ayuda a trabajar el
equilibrio, estirar espalda, desarrollar
fuerza en piernas. Mejora la concentración.
Masajea órganos abdominales.

\subsubsection*{Vrksāsana}

Vrksa = árbol.

\emph{Instrucciones.}
\begin{enumerate}
\item Comenzar en Tadāsana.
\item Levantar pierna derecha y llevar la
planta al muslo izquierdo. Si esto no es
posible puede ser al tobillo izquierdo.
\item Manos en la cadera y respiración
profunda.
\item Lentamente llevar las manos estiradas
por arriba de la cabeza y juntar las palmas.
\end{enumerate}

\emph{Acciones.}
\begin{itemize}
\item Estirar bien espalda, tracción de
columna. Estirar los costados.
\item Relajar el rostro y respirar profundo.
\item Pierna de abajo bien estirada.
\item Coxis adentro.
\item Abrir la cadera de la pierna de arriba.
\end{itemize}

\emph{Salida.}
\begin{enumerate}
\item Devolver brazos a los costados.
\item Pierna derecha abajo.
\item Hacer el otro lado.
\end{enumerate}

\emph{Beneficios.} Apertura de cadera,
estiramiento de espalda y equilibrio.
Concentración.

\subsubsection*{Utthita Hasta Pādānsuṣṭāsana}

Uttihita = extendido, estirado. Hasta = mano.
Pādānguṣṭa = dedo gordo del pie.

\emph{Variantes.} De frente sin cinturón, de
frente con cinturón y pared, de lado y con
torsión.

\emph{Instrucciones.}
\begin{enumerate}
\item Comenzar en Tadāsana.
\item Levantar el pie derecho y tomar el dedo
gordo del pie con la mano derecha. Mano
izquierda en cintura.
\item Estirar la pierna en un ángulo 45º
\end{enumerate}

\emph{Acciones.}
\begin{itemize}
\item Estirar columna y proyectar esternón.
\item Meter coxis.
\item Respirar profundo.
\end{itemize}

\emph{Salida.}
\begin{enumerate}
\item En una exhalación devolver pie al piso.
\item Hacer el otro lado.
\end{enumerate}

\emph{Adaptaciones.} Para facilitar la
postura se usa un cinturón en el pie. Así es
más fácil estirar la columna.

\emph{Beneficios.} Mejora el equilibrio, la
concentración, fortalece piernas y mejora la
flexibilidad de la cadera.

\medskip\noindent\emph{Grupo 6. Posturas en
flexión}

\subsubsection*{Uttānāsana}

Ut = intensidad, delibración. Tān = estirado,
extendido.

\emph{Instrucciones.}
\begin{enumerate}
\item Comenzar en Tadāsana con Ūrdhva
Hastāsana.
\item Llevar lentamente las manos al piso
estirando sobre todo la espalda baja. Toda la
columna está profundamente estirada duranete
el movimiento.
\item Finalmente las manos están en el piso,
si es posible toda la palma al suelo y la
frente hacia las rodillas
\end{enumerate}

\emph{Acciones.}
\begin{itemize}
\item Respirar profundo y en las exhalaciones
profundizar acercando más la frente a las
rodillas.
\item Para profundizar se puede tomar de los
tobillos y usar la fuerza de los brazos para
acercarse más.
\item Hombros lejos de orejas.
\end{itemize}

\emph{Salida.}
\begin{enumerate}
\item Elevar el torso lentamente percibiendo
cómo se reacomoda cada vértebra hasta llegar
a Tadāsana con Ūrdhva Hastāsana.
\end{enumerate}

\emph{Contraindicaciones.} Realizar sólo
hasta su punto medio (Ardha Uttanāsana)
durante el embarazo o si se presentan dolores
de columna.

\subsubsection*{Pādānguṣṭāsana}

Pādānguṣṭa = dedo gordo del pie.

\emph{Instrucciones.}
\begin{enumerate}
\item Comenzar en Tadāsana con Ūrdhva
Hastāsana..
\item Llevar lentamente las manos al piso
estirando sobre todo la espalda baja. Toda la
columna está profundamente estirada duranete
el movimiento.
\item Finalmente los dedos índice y medio
toman el dedo gordo de cada pie. La columna
sigue estirada y la mirada al frente.
\item Por último llevamos frente a rodillas
flexionando la columna.
\end{enumerate}

\emph{Acciones.}
\begin{itemize}
\item Respirar profundo y en las exhalaciones
profundizar acercando más la frente a las
rodillas.
\item Hombros lejos de orejas.
\end{itemize}

\emph{Salida.}
\begin{enumerate}
\item Inhalando, lentamente regresar a
Tadāsana.
\end{enumerate}

\emph{Beneficios.} Masaje en órganos
abdominales.

\emph{Contraindicaciones.} Realizar sólo
hasta su punto medio durante el embarazo o si
se presentan dolores de columna.

\subsubsection*{Pādāhastāsana}

Pādā = pie. Hasta = mano

\emph{Instrucciones.}
\begin{enumerate}
\item Comenzar en Tadāsana con Ūrdhva
Hastāsana..
\item Llevar lentamente las manos al piso
estirando sobre todo la espalda baja. Toda la
columna está profundamente estirada duranete
el movimiento.
\item Las palmas de las manos entran en
contacto con las plantas de los pies. Dorsos
de las manos al piso. La columna sigue
estirada y la mirada al frente.
\item Por último llevamos frente a rodillas
flexionando la columna.
\end{enumerate}

\emph{Acciones.}
\begin{itemize}
\item Respirar profundo y en las exhalaciones
profundizar acercando más la frente a las
rodillas.
\item Hombros lejos de orejas.
\end{itemize}

\emph{Salida.}
\begin{enumerate}
\item Inhalando, lentamente regresar a
Tadāsana.
\end{enumerate}

\emph{Beneficios.} Masaje en órganos
abdominales.

\emph{Contraindicaciones.} Realizar sólo
hasta su punto medio durante el embarazo o si
se presentan dolores de columna.

\subsubsection*{Prasārita Pādottānāsana}

Prasārita = expandido, extendido. Pāda = pie.

\emph{Instrucciones.}
\begin{enumerate}
\item Comenzar en Tadāsana.
\item Separar las piernas dando un brinco o
un pequeño salto aproximadamente a un metro
de distancia.
\item Estirar ambos brazos paralelos al piso
con las palmas mirando hacia abajo.
\item Estirar bien la espalda separando
vértebras y abriendo el pecho, y llevar el
pecho a una flexión hacia en frente llevando
la cabeza abajo de las caderas.
\item Las manos pueden ir en el suelo con los
codos flexionados o tomando los tobillos.
\end{enumerate}

\emph{Acciones.}
\begin{itemize}
\item Estirar completamente la columna.
\item Relajar el cuello.
\item Activar rótulas y muslos activos.
\item Respiración profunda.
\item Se pueden hacer variantes con los
brazos: antes de la flexión de columna hacer
manos de Paschima Namaskarāsana o entrelazar
las manos por atrás del cuerpo estirando
brazos.
\end{itemize}

\emph{Salida.}
\begin{enumerate}
\item Microflexionar rodillas y llevar manos
a rodillas para reincorporarse lentamente.
\item Juntar piernas y terminar en Tadāsana.
\end{enumerate}

\emph{Beneficios.} Estirar espalda,
fortalecer el cuello, mejora circulación,
tonifica abdomen y mejora la respiración.

\emph{Contraindicaciones.} No realizar
durante el embarazo o si se presentan
lesiones de espalda o dolor de cabeza.

\subsubsection*{Parighāsana}

Parighā = es el travesaño de una puerta.

\emph{Instrucciones.}
\begin{enumerate}
\item Comenzar hincado con las rodillas en el
piso y los metatarsos de los dedos del pie
apoyados.
\item Estirar el pie derecho hacia el lado
derecho abriendo ligeramente el pie (rotación
externa). La rodilla queda bien estirada.
\item Extender los brazos horizontalmente con
las palmas mirando hacia abajo.
\item Exhalar y hacer flexión lateral de
columna llevando la mano derecha al tobillo
derecho. Estirar el brazo izquierdo hacia el
lado derecho haciendo contacto con la oreja
izquierda.
\item Respirar profundo. Si se quiere entrar
a la postura completa, profundizar en la
flexión lateral hasta llevar ambas manos al
pie derecho.
\end{enumerate}

\emph{Acciones.}
\begin{itemize}
\item Respirar profundo y en las exhalaciones
profundizar en la flexión.
\item Meter abdomen.
\item Estirar bien la axila izquierda.
\item Relajar músculos del cuello.
\item Conectar con la compactación en los
órganos abdominales.
\end{itemize}

\emph{Salida.}
\begin{enumerate}
\item En un a inhalación devolver el tronco a
su posición normal con la columna estirada.
\item Flexionar la rodilla derecha y devolver
la pierna a la posición hincado.
\item Hacer el otro lado.
\end{enumerate}

\emph{Beneficios.} Estiramiento en la pelvis.
Mejora los órganos abdominales. Revitaliza
espalda.

\emph{Contraindicaciones.} No realizar
durante el embarazo o si se presentan
lesiones de espalda o dolor de cabeza.

\subsubsection*{Pārśvottānāsana}

\emph{Instrucciones.}
\begin{enumerate}
\item Entrar a Parśva Hasta Pādāsana.
\item Estirar profundamente la espalda desde
lumbares hasta el cuello y flexionar la
columna y elevar los brazos verticalmente
\item En una exhalación llevar frente a
rodilla. Las manos tocan el piso.
\end{enumerate}

\emph{Acciones.}
\begin{itemize}
\item Estirar la espalda desde lumbares.
\item Conectar con la compactación del
abdomen.
\item Rótulas activas y piernas bien
estiradas.
\item Presionar el borde externo del pie de
adelante al piso.
\item Las variantes pueden ser con las palmas
de las manos en contacto o en manos de
Paschima Namaskarāsana.
\end{itemize}

\emph{Salida.}
\begin{enumerate}
\item Vértebra por vértebra devolver el torso
a su posición natural con la espalda y los
brazos bien estirados.
\item Regresar los brazos a la horizonatal
para volver a Parśva Hasta Pādāsana y hacer
el otro lado.
\end{enumerate}

\emph{Beneficios.} Estira columna, fortalece
piernas, estimula órganos abdominales, mejora
el equilibrio y relaja el cerebro.

\subsubsection*{Ūrdhva Prasārita Eka Pādāsana}

Ūrdhva = arriba, alto. Prasārita = estirado,
extendido. Eka = uno. Pāda = pie.

\emph{Instrucciones.}
\begin{enumerate}
\item Comenzar en Vīrabhadrāsana II.
\item Llevar el brazo izquierdo hacia
adelante hasta alcanzar el brazo derecho.
\item Continuar con esta trayectoria
flexionando el torso y llevando ambas manos
al piso y levantando el pie de atrás.
\item Tomar con la mano derecha el tobillo
derecho y llevar frente a rodilla.
\end{enumerate}

\emph{Acciones.}
\begin{itemize}
\item Respirar profundo y en las exhalaciones
acercar más frente a rodilla.
\item Estirar bien la pierna de atrás y
llevar el pie al techo.
\item Activar rótula derecha y estirar bien
la pierna derecha.
\item Relajar cuello.
\item Respirar bien profundo y relajar la
mente.
\end{itemize}

\emph{Salida.}
\begin{enumerate}
\item Llevar ambas manos al piso.
\item Simultáneamente llevar ambos brazos a
la horizontal, levantar el torso y llevar la
pierna izquierda al piso.
\item Girar las piernas para entrar a
Vīrabhadrāsana II del lado contrario y hacer
el otro lado.
\end{enumerate}

\emph{Beneficios.}
\begin{itemize}
\item Estimula sistema digestivo.
\item Fortalece piernas.
\item Estimula sistema nervioso.
\end{itemize}

\emph{Contraindicaciones.} Lesiones en
tobillos o lumbares.

\subsection{Posturas sentados}
\label{section-posturas-sentados}

\subsubsection*{Daṇḍāsana}

Daṇḍa = apoyo, bastón, vara.

\emph{Instrucciones.}
\begin{enumerate}
\item Sentarse en el piso con las piernas
bien estiradas, rótulas activas y pies en
flex.
\item El torso está recto con la columna
estirada hacia el techo.
\item Las manos a un costado de las caderas
con las palmas en el piso.
\end{enumerate}

\emph{Acciones.}
\begin{itemize}
\item Mantener brazos, piernas y espalda bien
estiradas.
\item Relajar hombros, cuello y rostro.
\item Pies bien en flex.
\item Estómago adentro.
\item Pecho abierto.
\end{itemize}

\emph{Beneficios.} Calma la mente, refuerza
espalda y piernas. Es una postura de
transición.

\subsubsection*{Sukhāsana}

Sukha = plenitud, amplitud.

\emph{Instrucciones.}
\begin{enumerate}
\item Sentarse con las piernas cruzadas,
isquines en el suelo sobre tantas cobijas
como sean necesarias para que las rodillas no
estén por arriba de las caderas.
\end{enumerate}

\emph{Acciones.}
\begin{itemize}
\item Estirar espalda separando vértebras y
llevando coronilla al techo.
\item Estirar el cuello.
\item Relajar rostro.
\item Los brazos pueden descansar sobre
rodillas o muslos con las palmas hacia
arriba, abajo o en algún Mudrā.
\item Brazos estirados o codos ligeramente
flexionados.
\item Respirar bien profundo y concentrarse.
\end{itemize}

\emph{Beneficios.} Estira la espalda, aumenta
la concentración, relajación. Es una
excelente postura para meditar y hacer
pranayama.

\emph{Contraindicaciones.} Sólo es importante
cuidar la altura de las rodillas usando los
soportes que sean necesarios.

\subsubsection*{Baddha Koṇāsana}

Baddha = ligado, atado. Koṇa = ángulo.

\emph{Instrucciones.}
\begin{enumerate}
\item Entrar a Daṇḍāsana.
\item Flexionar las rodillas y hacer entrar
en contacto las palmas de los pies.
\item Entrelazar las manos alrededor de los
pies.
\end{enumerate}

\emph{Acciones.}
\begin{itemize}
\item Estirar la espalda y los brazos.
\item Estirar el cuello pero relajar los
hombros.
\item Relajar el rostro.
\item Conectar con la apertura de cadera.
Relajar la cadera para aumentar la
flexibilidad.
\end{itemize}

\emph{Salida.}
\begin{itemize}
\item Soltar el entrelazado de las manos.
\item Estirar las piernas y mover ligeramente
para devolver la circulación.
\end{itemize}

\emph{Beneficios.} Apertura de caderas,
concentración, estiramiento de espalda,
relajación.

\emph{Contraindicaciones.} Si hay tensión en
rodillas o en caderas se puede hacer la
postura restaurativa.

\subsubsection*{Siddhāsana}

Siddha = son seres semi-divinos, puros y
sagrados que poseen poderes sobrenaturales
llamados siddhis.

\emph{Instrucciones.}
\begin{enumerate}
\item Entrar a Daṇḍāsana.
\item Flexionar la pierna izquierda llevando
al pie al coxis.
\item Flexionar la pierna derecha y llevar el
pie derecho al coxis. El pie derecho debe ir
sobre el tobillo izquierdo, de forma que la
planta del pie hace contacto con el muslo y
el talón con el pubis.
\item Estirar los brazos y descansar los
dorsos de las manos sobre las rodillas.
\item Para salir simplemente desdoblar las
piernas y volver a Daṇḍāsana.
\end{enumerate}

\emph{Acciones.}
\begin{itemize}
\item Estirar bien espalda.
\item Estirar músculos del cuello y relajar
hombros.
\item Relajar rostro y respirar profundo.
\end{itemize}

\emph{Beneficios.} Ayuda a relajar cóccix,
flexibilidad en cadera, rodillas y talones.
Es sumamente relajante y buena para meditar.

\subsubsection*{Padmāsana}

Padma = loto

\emph{Instrucciones.}
\begin{enumerate}
\item Entrar a Daṇḍāsana.
\item Flexionar la rodilla izquierda un
ángulo que sea perfectamente cómodo y
relajado.
\item Tomar el pie derecho con ambas manos y
girar la pierna unas cuantas veces para
estimular la articulación de la cadera (de
cualquier forma se recomienda hacer esta
postura después de algunas otras que sean
preparativas para cadera).
\item Llevar el pie derecho sobre el muslo
izquierdo lo más cerca de la ingle posible. Y
flexionar la rodilla izquierda para que el
pie izquierdao esté por debajo de la rodilla
derecha. (Ardha Padmāsana).
\item Respirar unos segundos.
\item Tomar el pie izquierdo con ambas manos
y llevarlo sobre el muslo derecho lo más
cerca que se pueda a la ingle derecha.
\item Estirar los brazos sobre las rodillas
manteniendo la espalda y el cuello bien
estirados.
\item Hacer las variantes pertinentes.
\item Para salir liberar primero la pierna
izquierda y luego la derecha. Hacer pequeños
movimientos con las piernas estiradas para
recuperar la circulación. Hacer el otro lado.
\end{enumerate}

\emph{Acciones.}
\begin{itemize}
\item Estirar totalmente la espalda separando
cada vértebra desde sacro hasta el cuello.
\item Relajar hombros y cuello.
\item Relajar rostro.
\item Relajar la mente.
\item Coronilla al cielo.
\item Se pueden explorar diferentes Mudrās
con las manos.
\end{itemize}

\emph{Beneficios.} Es una excelente postura
para meditar y relajarse. Ayuda con
flexibilidad de cadera, rodillas y tobillos.
Ayuda con la circulación y tonifica abdomen.

\subsubsection*{Upaviṣṭha Koṇāsana}

Upaviṣṭha = sentado. Koṇa = ángulo.

\emph{Variantes.} Brazos al piso y adelante,
o tomando dedos gordos de los pies.

\emph{Instrucciones.}
\begin{enumerate}
\item Entrar a Daṇḍāsana.
\item Separar las piernas simultáneamente y
comprobar que estén en el mismo ángulo.
\item Llevar el torso hacia adelante
estirando profundamente las lumbares y
flexionando sólo hasta el final manteniendo
la intención de llevar abdomen al piso. Los
brazos estirados hacia adelante.
\item Para salir devolver la espalda a la
vertidal y volver a juntar piernas para
regresar a Daṇḍāsana.
\end{enumerate}

\emph{Acciones.}
\begin{itemize}
\item Activar rótulas y estirar bien las
piernas.
\item Flexionar pies.
\item Relajar cabeza.
\item Concentrarse en espalda baja.
\item Estirar los brazos hacia adelante.
\item En las exhalaciones profundizar en la
postura.
\item Una variante que ayuda a profundizar es
tomando los dedos gordos de los pies con los
dedos índice y medio.
\end{itemize}

\emph{Beneficios.} Fortalece columna
vertebral, abre caderas, fortalece piernas.

\emph{Contraindicaciones y adaptaciones.}
Sentarse sobre una manta puede ayudar a
profundizar en la flexion. Durante el
embarazo no flexionar la espalda sino bajar
lo que se posible con la espalda
completamente recta.

\subsubsection*{Vajrāsana}

Vajra = rayo, es el arma de Indra.

\emph{Instrucciones.}
\begin{enumerate}
\item Comenzar hincado con las rodillas
juntas y los empeines al piso.
\item Sentarse sobre los talones. Y llevar
los brazos bien estirados a rodillas. Espalda
estirada.
\item Para salir soltar los brazos y volver a
posición hincado.
\end{enumerate}

\emph{Acciones.} Estirar bien columna,
relajar cuello y rostro.

\emph{Beneficios.} Es sumamente relajante.
Ayuda con flexibilidad en rodillas y sirve
para meditar.

\emph{Contraindicaciones y adaptaciones.} Si
hay tensión en empeines o rodillas se puede
usar un soporte en espalda baja.

\subsubsection*{Vīrāsana}

Vīra = héroe, guerrero.

\emph{Instrucciones.}
\begin{enumerate}
\item Comenzar hincado con las rodillas
juntas y los empeines al piso.
\item Separar los pies para que los isquiones
se apoyen en el piso. Las palmas van sobre
las rodillas con los codos estirados.
\item Para salir soltar los brazos y volver a
posición hincado.
\end{enumerate}

\emph{Acciones.} Estirar bien columna,
relajar cuello y rostro.

\emph{Beneficios.} Es sumamente relajante.
Ayuda con flexibilidad en rodillas y sirve
para meditar.

\emph{Contraindicaciones y adaptaciones.} Si
hay tensión en empeines o rodillas se puede
usar un soporte en espalda baja. Para entrar
puede ser conveniente apoyar la frente o la
coronilla al piso para hacer una ligera
rotación externa de los gemelos
(pantorrillas). Esto facilita la entrada.

\subsubsection*{Parvattāsana}

Parvata=montaña

\emph{Instrucciones.}
\begin{enumerate}
\item Entrar a Padmāsana.
\item Entrelazar los dedos de las manos y
elevar los brazos sobre la cabeza. Las palmas
miran el cielo.
\item Bajar los brazos, cambiar el
entrelazado de los dedos y volver a elevar
brazos.
\end{enumerate}

\emph{Acciones.} Estirar bien columna,
relajar cuello y rostro.

\emph{Beneficios.} Además dáe los beneficios
de Padmāsana, relaja hombros y estira
músculos y órganos abdominales.

\emph{Otras variantes de Padmāsana.} Flexión
al frente, torsión, Tolāsana (Tola = balance,
equilibrio) y recostado.

\subsubsection*{Gomukhāsana}

Go = vaca. Mukha = cara. Gomukha es aquel
cuya cara parece de vaca, y también es el
nombre de un instrumento musical.

\emph{Instrucciones.}
\begin{enumerate}
\item Comenzar en Daṇḍāsana.
\item Elevar la cadera del suelo apoyándose
bien en las manos y flexionar la rodilla
izquieda hasta sentarse en el talón
izquierdo.
\item Ayudándose con las manos flexionar la
piernda derecha y cruzarla por encima de la
izquierda.
\item Acercar los pies hasta que los tobillos
se toquen (si es posible).
\item Estirar los brazos en una línea
horizontal al piso. La palma izquierda mira
al frente y la derecha atrás.
\item Llevar el izquierdo por detrás de la
cabeza y el dereho por detrás de la espalda
baja.
\item Entrelazar las manos con todos los
dedos juntos por detrás de la espalda. Ésta
es la postura completa.
\item Para salir, soltar el entrelazado de
manos y regresar lentamente los brazos
paralelos al piso.
\item Regresar las manos al piso y ayudarse
de ellas para deshacer el entrelazado de las
piernas.
\item Estirar y relajar las piernas. Hacer el
otro lado.
\end{enumerate}

\emph{Acciones.}
\begin{itemize}
\item Estirar bien la espalda, coronilla al
cielo.
\item Acercar más los tobillos entre sí.
\item Respirar con el diafragma y relajarse.
\item Al principio puede ser difícil lograr
el entrelazado de las manos: no dudar en usar
un cinturón.
\item Después de unas respiraciones en la
postura las articulaciones se relajan y es
posible ir más profundo. Observar esto y en
las exhalaciones acercar más las manos.
\item Hacer círculos milimétricos con el
hombro del brazo que está arriba ayuda a
desarrollar la flexibilidad. Este movimiento
puede ser incómodo: hay que encontrar el
punto en el que no hay dolor ni tensión (es
una buena oportunidad para practicar nuestra
habilidad de realizar movimientos o posturas
incómodas cómodamente: el truco no pasar
nuestro límite).
\item Los ``brazos de Gomukhāsana'' aparecen
en muchas otras posturas, como mencionamos
arriba en una variante de Tadāsana. Otra
variante, parte del Protocolo de Cintura
Escapular, consiste en colocar sólo el brazo
de abajo como en Gomukhāsana y recostarse
boca arriba sobre él, con la palma mirando al
piso; esto es sumamente relajante y poderoso
para flexibilizar esta articulación.
\end{itemize}

\emph{Beneficios.}
\begin{itemize}
\item Abre el pecho. Es excelente para
desarrollar flexibilidad en hombros y cintura
escapular.
\item Se trabajan fuertemente los músculos
deltoides y dorsal ancho.
\item Fortalece columna.
\item Flexibiliza rodillas y cadera.
\end{itemize}

\subsection{Posturas con torsión sentados}
\label{section-posturas-con-torsion-sentados}

\subsubsection*{Marīchyāsana}

Marīchi fue un sabio hijo de Brahmā y abuelo
de Sūrya (Dios del Sol).

\subsubsection*{Marīchyāsana I}

\emph{Instrucciones.}
\begin{enumerate}
\item Entrar a Daṇḍāsana.
\item Flexionar rodilla izquierda y acercar
el talón al perineo de tal forma que la
espinilla esté totalmente vertical. El
costado interno del pie está en contacto con
el muslo derecho.
\item Le pierna derecha se mantiene estirada
y el pie está en flex (o en punta, de acuerdo
a Iyengar).
\item Estirar bien la espalda con la columna
completamente vertical. Estirar el brazo
izquierdo verticalmente y luego llevarlo, sin
doblar el codo hacia en frente hasta que
quede paralelo al piso con la palma de la
mano mirando hacia el lado derecho.
\item Rotar el brazo en dirección de las
manecilla del reloj hasta que la palma mire
hacia el lado izquierdo.
\item Lentamente, con conciencia y control,
abrazando la espinilla izquierda llevar el
brazo izquierdo hacia la espalda hasta que el
dorso de la mano esté cerca el sacro.
\item Llevar el brazo por atrás del cuerpo
hasta alcanzar la mano izquierda por detrás
de la espalda. Entrelazar las manos.
\item Respirar profundo y estirar bien la
espalda verticalmente haciendo tracción de
columna.
\item En una exhalación llevar frente a
rodillas para hacer una flexión profunda.
\item Para salir, liberar el entrelazado y
estirar la espalda vértebra por vértebra.
Hacer el otro lado.
\end{enumerate}

\emph{Acciones.}
\begin{itemize}
\item En las exhalaciones acercar la frente a
rodilla.
\item Si no es posible entrelazar las manos,
se puede usar cinturón.
\item Al principio puede resultar una postura
incómoda. Explorar pacientemente cuál es la
flexibilidad que se tiene para ir tan
profundo como sea posible sin estresar las
articulaciones. Entrar a la postura aunque
sea superficialmente es bueno para entender
cómo funciona y con el tiempo desarrollar la
flexibilidad para profundizar y obtener sus
beneficios.
\end{itemize}

\emph{Beneficios.}
\begin{itemize}
\item Los órganos abdominales se contraen
fuertemente, aumentando la circulación y
rejuveneciendo los tejidos.
\item Se ejercitan la espalda alta y media.
\item Junto con Jānu Sīrsāsana, Triang
Mukhaipāda Paschimottānāsana, Ardha Baddha
Padma Paschimottānāsana, estas cuatro
posturas nos ayudan a desarrollar la
flexibilidad en piernas y espalda para entrar
correctamente a Paschimottānāsana
\end{itemize}

\subsubsection*{Marīchyāsana II}

\emph{Instrucciones.}
\begin{enumerate}
\item Comenzar en Daṇḍāsana.
\item Flexionar la pierna izquierda en
Vīrāsana y dejar la derehca estirada.
\item Realizar el mismo procedimiento de
entrelazado que en Marīchyāsana I: elevar el
brazo derecho haciendo tracción de columna,
llevarlo paralelo al piso, rotarlo en
dirección antihoraria para que la palma mire
al lado derecho con el pulgar abajo, y
flexionarlo abrazando la espinilla derecha
hasta entrelazar las manos detrás de la
espalda.
\item Respirar estirando la columna hacia
arriba y en una exhalación llevar frente a
rodilla en una profunda flexión. Respirar
profundo.
\item Para salir, liberar el entrelazado y
estirar la espalda vértebra por vértebra.
Hacer el otro lado.
\end{enumerate}

\emph{Acciones.} Las mismas que en
Marīchyāsana I.

\emph{Variante.} Esta postura también se
puede hacer con el pie izquierdo en Ardha
Padmāsana. (Marīchyāsana IV es la versión en
torsión de esta variante).

\subsubsection*{Marīchyāsana III}

\emph{Instrucciones.}
\begin{enumerate}
\item Entrar a Daṇḍāsana.
\item Flexionar rodilla izquierda y acercar
el talón al perineo de tal forma que la
espinilla esté totalmente vertical. El
costado interno del pie está en contacto con
el muslo derecho.
\item Le pierna derecha se mantiene estirada
y el pie está en flex (o en punta, de acuerdo
a Iyengar).
\item Estirar bien la espalda con la columna
completamente vertical. Apoyar la mano
izquierda en el piso detrás de la espalda y
elevar el brazo derecho verticalmente.
\item Girar 90º a la izquierda sin doblar la
espalda hasta que el hombro esté más allá de
la pierna izquierda, y llevar el antebrazo
derecho sobre el muslo.
\item Si es posible, estirar el brazo derecho
para tomar la parte externa de la espinilla
derecha.
\item Si el paso anterior es fácil, flexionar
el codo derecho abrazando la pierna izquierda
y llevar la mano en dirección de la espalda
baja. Las manos se entrelazan (o bien la
derecha toma la muñeca izquierda) por detrás,
abrazando a la pierna flexionada.
\item Para salir primero liberar las manos y
regresar los brazos al paso anterior.
Lentamente salir de la torsión y hacer el
otro lado.
\end{enumerate}

\emph{Acciones.}
\begin{itemize}
\item Estirar bien la espalda para que la
columna continúe siendo una línea vertical
perpendicular al piso.
\item En las exhalaciones profundizar en la
torsión llevando la mirada hacia la pared de
atrás y fortaleciendo el soporte de las manos
o brazos.
\item Activar bien la pierna de abajo.
\end{itemize}

\emph{Beneficios.}
\begin{itemize}
\item Incrementa flexibilidad en la columna,
mejora la circulación.
\item Estimula fuertemente los órganos
abdominales: intestinos, hígado, bazo.
\item Fortalece músculos del cuello.
\item Mejora la articulación del hombro, es
un fuerte trabajo en el deltoides.
\end{itemize}

\emph{Contraindicaciones.} Iyengar recomienta
a personas de complexión delgada hacer otras
posturas, pues de ésta no se beneficiarán
tanto.

\subsubsection*{Marīchyāsana IV}

\emph{Instrucciones.}
\begin{enumerate}
\item Entrar a Daṇḍāsana.
\item Flexionar la pierna derecha y llevar el
dorso del pie sobre el muslo izquierdo en
Ardha Padmāsana.
\item Flexionar rodilla izquierda y acercar
el talón al perineo de tal forma que la
espinilla esté totalmente vertical. El pie
derecho queda entre el muslo el abdomen,
apretando el ombligo. Estirar el brazo
izquierdo apoyando la palma en el piso.
\item Inhalar y estirar la columna en
tracción, y al exhalar girar el cuerpo 90º a
la izquierda llevando el antebrazo derecho
sobre el muslo izquierdo.
\item Si es posible, llevar la axila y luego
todo el hombro por delante del muslo
izquierdo.
\item Si lo anterior es fácil, estirar el
brazo derecho, respirar, y luego flexionar el
codo llevando la mano al sacro sin liberar la
espinilla izquierda. La mano izquierda se
encuentra con la derecha y se entrelazan (o
mano derecha toma muñeca izquierda).
\item Salir lentamente de la torsión y hacer
el otro lado.
\end{enumerate}

\emph{Acciones.} Mismas que en Marīchyāsana
III.

\emph{Variante.} Esta postura también se
puede hacer con el pie derecho en Vīrāsana
(Marīchyāsana II es la versión en flexión de
esta variante).

\emph{Beneficios.} Además de los beneficios
de Marīchyāsana III, la presión del pie
contra el ombligo rejuevenece los nervios en
esta zona.

\subsubsection*{Bharadvājāsana}

Bharadvāja era el padre de Drona, el maestro
militar de los Kauravas y los Pāṇḍavas, que
lucharon en la guerra descrita en el
Mahābhārata.

\subsubsection*{Bharadvājāsana I}

\emph{Instrucciones.}
\begin{enumerate}
\item Entrar a Daṇḍāsana.
\item Flexionar las rodillas como en
Vajrāsana y sentarse en el piso dejando los
dos pies del lado izquierdo.
\item Girar el torso 45º hacia el lado
derecho manteniendo la columna bien estirada.
\item Estirar el brazo izquierdo y apoyar el
dorso de la mano en el lado externo del muslo
izquierdo. Si es posible, la palma de la mano
toca el piso y el dorso de la mano está en
contacto con la rodilla.
\item Flexionar el codo derecho y tratar de
alcanzar el bícep izquierdo. En la postura
completa hay contacto.
\item Llevar la mirada hacia la derecha por
ariiba del hombro.
\item Para salir de la postura, regresar
primero cabeza, luego brazos y torso, y al
final estirar piernas. Relajar las piernas y
hacer el otro lado.
\end{enumerate}

\emph{Acciones.}
\begin{itemize}
\item Estirar la columna en tracción
separando vértebra por vértebra.
\item En las exhalaciones profundizar en la
torsión llevando todo el torso hacia atrás,
guiándose con la mirada.
\item Relajar concientemente la articulación
del hombro para desarrollar la flexibilidad
necesaria para que el brazo de atrás alcance
al otro brazo.
\end{itemize}

\emph{Beneficios.} Las dos variantes de
Bharadvājāsana son posturas accesibles para
quienes no tienen mucha flexibilidad (sobre
todo la variante I). Son buenas para entender
en qué consisten las torsiones y
familiriazarse con la sensación.

Mejoran la flexibilidad de la columna, sobre
todo en las regiones lumbar y dorsal.
Recomendadas para personas con artitritis.

\subsubsection*{Bharadvājāsana II}

\emph{Instrucciones.}
\begin{enumerate}
\item Entrar a Daṇḍāsana.
\item Flexionar la pierna derecha en Vīrāsana
y dejar la izquierda estirada.
\item Flexionar suavemente la pierna
izquierda y llevarla sobre el muslo derecho
como en Ardha Padmāsana
\item Girar el torso 45º hacia el lado
derecho manteniendo la columna bien estirada.
\item Estirar el brazo izquierdo y apoyar el
dorso de la mano en el lado externo del muslo
izquierdo. Si es posible, la palma de la mano
toca el piso y el dorso de la mano está en
contacto con la rodilla.
\item Flexionar el codo derecho y tratar de
alcanzar el bícep izquierdo. En la postura
completa hay contacto.
\item Llevar la mirada hacia la derecha por
ariiba del hombro.
\item Para salir de la postura, regresar
primero cabeza, luego brazos y torso, y al
final estirar piernas. Relajar las piernas y
hacer el otro lado.
\end{enumerate}

\emph{Acciones.} Mismas que en Bharadvājāsana
I. Si no es posible tomar el pie con el brazo
de atrás, apoyarlo en el piso.

\emph{Beneficios.} Ver Bharadvājāsana I.

\subsubsection*{Ardha Matsyendrāsana}

Ardha=mitad. Cuenta la historia que el Señor
Śiva fue a una isla desierta con su consorte
Pārvati para explicarle los misterios del
Yoga. Un pez cerca de la orilla se acercó
sigilosamente para escuchar. Sabiendo que el
pez había aprendido el Yoga, Śiva le echó
agua encima, causando que el pez adoptara su
forma divina para convertirse en Matsyendra
(Señor de los Peces). Paripūrna
Matsyendrāsana está dedicada a Matsyendra, y
Ardha Matsyendrāsana es una versión más
sencilla).

\emph{Instrucciones.}
\begin{enumerate}
\item Entrar a Daṇḍāsana.
\item Flexionar la rodilla izquierda y apoyar
empenie sobre el piso para sentarse sobre el
talón izquierdo. Acomodar como sea necesario
el cuerpo para encontrar comodidad en la
postura.
\item Flexionar la pierna derecha y llevar el
pie al piso por fuera del muslo izquierdo.
\item Estirar bien la columna y girar el
torso 90º hacia la derecha hasta que la axila
y el hombro estén más allá de la rodilla
derecha. Apoyar el antebrazo izquierdo sobre
el muslo derecho. El brazo derecho descansa
sobre el muslo izquierdo. Ésta puede ser la
postura final.
\item Si es fácil entrar al paso anterior,
estirar el brazo izquierdo para tomar la
rodilla izquierda. Respirar profundo.
\item Si incluso lo anterior resulta
sencillo, estirar el brazo izquierdo paralelo
al piso y rotarlo para que la palma de la
mano esté hacia la izquierda con el pulgar
abajo.
\item Flexionar el codo izquierdo abrazando
la pierna como en Marīchyāsana y entrelazar
la manos por detrás del sacro. Mantener la
columna bien estirada.
\item Para salir, liberar primero las manos y
regresar el pecho de la torsión. Liberar la
pierna derecha y por último la izquierda.
Relajar un poco las piernas estiradas y hacer
el otro lado.
\end{enumerate}

\emph{Acciones.}
\begin{itemize}
\item Si es complicado el apoyo sobre el
tobillo se puede experimentar la postura
sentado en el piso. Ésta es una torsión muy
accesible que se puede practicar como postura
de transición en una secuencia.
\item También es buena opción estirar el
brazo de atrás para apoyar la mano en el piso
y ayudarse de ella para profundizar en la
torsión.
\item Cada vez que se exhala, llevando más la
mirada hacia atrás, manteniendo la espalda
bien recta y estirada.
\item Relajar bien el cuello y el rostro.
\end{itemize}

\emph{Beneficios.}
\begin{itemize}
\item La postura completa es más profunda que
Marīchyāsana III, y los beneficios son los
mismos pero en mayor proporción.
\item Se estimula el abdomen bajo, ayudando a
mantener sana la vejiga y la próstata en
hombres.
\end{itemize}

\subsubsection*{Mālasāsana I, II}

Mālā = guirnalda.

\emph{Instrucciones.}
\begin{enumerate}
\item Comenzar en cuclillas. Si esto no es
sencillo, se puede
\begin{itemize}
\item Separar más los pies.
\item Abrir el ángulo de los pies con las
puntas hacia afuera.
\item Usar un soporte (cobija o mat) sobre
los talones.
\end{itemize}
\item Separar las rodillas y llevar torso
quepa entre ellas.
\item Apoyar los antebrazos en las rodillas
\item Mālasāna I:
\begin{enumerate}
\item Juntar las palmas en Namaskarāsana.
Ésta puede ser la postura final.
\item Si se tiene buena flexibilidad, estirar
los brazos y apoyar las palmas al piso
manteniendo la columna bien estirada.
\item Flexionar ambos brazos y entrelazar las
manos a la altura del sacro manteniendo la
espalda recta.
\item Exhalar y flexionar la columna llevando
frente al piso. Ésta es la postura completa.
\item Para salir, primero regresar de la
torsión estirando la columna y luego liberar
los brazos. Apoyar las manos en el piso para
relajar piernas.
\end{enumerate}
\item Mālāsana II:
\begin{enumerate}
\item Flexionar los brazos y tomar las
espinillas por atrás manteniendo la columna
bien estirada.
\item Exhalar y flexionar la columna llevando
frente al piso. En la postura completa hay
contacto.
\item Para salir, regresar primero de la
flexión estirando bien la espalda, luego
librar los brazos y por último las piernas.
\end{enumerate}
\end{enumerate}

\emph{Acciones.}
\begin{itemize}
\item La postura se puede quedar en fase de
tracción de columna o se puede entrar en la
flexión. En cada caso traer la mente a la
columna y en las exhalación estirar (tracción
de columna) o llevar más la frente al piso
(flexión de columna).
\item Relajar la cadera y las ingles para que
en las exhalaciones bajen los isquiones al
piso sin hacer esfuerzo.
\item Separar rodillas entre sí y llevar más
el pecho al frente sin doblar la columna (a
menos que se esté en la flexión).
\end{itemize}

\emph{Beneficios.}
\begin{itemize}
\item Se estimulan los órganos abdominales y
se deasarrolla flexibilidad en columna.
\item Se fortalecen las piernas y se abre
cadera.
\item Mālasāna II ayuda a las mujeres que
sufren de dolores de espalda durante la
menstruación.
\end{itemize}

\subsection{Posturas de flexión}
\label{section-posturas-de-flexion}

\subsubsection*{Jānu Sīrṣāsana}

Jānu = rodilla. Sīrṣa = cabeza.

\emph{Instrucciones.}
\begin{enumerate}
\item Entrar a Daṇḍāsana.
\item Flexionar la rodilla izquierda y llevar
la planta del pie izquierdo en contacto con
el costado interno del muslo derecho.
\item Estirar ambos brazos verticalmente,
haciendo tracción de columna y llevando
coronilla al cielo, y rotar hacia la pierna
derecha.
\item Llevar los brazos bien estirados al pie
derecho y tomar la planta del pie. Si esto no
es posible se puede tomar muslo, espinilla o
dedo gordo del pie. Mantener bien estirarda
la columna abriendo el pecho. Tomar unas
cuantas respiraciones.
\item En una exhalación entrar a una flexión
profunda desde las vértebras lumbares
llevando abdomen a muslo, esternón a rodilla
y barbilla a espinilla.
\item Si el paso anterior es fácil, llevar el
cuerpo al piso por cualquiera de los dos
lados de la pierna. En este caso, no permitir
que la pierna gire. Ésta es la variante más
profunda.
\item Para salir, primero estirar la columna
sin liberar las manos como en el paso 4.
Luego soltar el agarre y estirar ambas
piernas. Hacer el otro lado.
\end{enumerate}

\emph{Acciones.}
\begin{itemize}
\item Llevar la rodilla flexionada hacia
atrás para abrir la cadera. En la postura
completa la rodilla está por detrás del
cuerpo, y el ángulo entre las piernas es
mayor a 90º.
\item En cada exhalación estirar bien la
espalda desde las vértebras lumbares hasta
las dorsales. Al principio es posible que el
torso esté lejos de la pierna. Es importante
no concentrarse en disminuir esta distancia y
estirar bien la espalda desde las vértebras
lumbras, sin encorvarse. Una y otra vez
recordar esta acción.
\item Es posible hacer esta postura
reteniendo el aire después de las
exhalaciones.
\end{itemize}

\emph{Beneficios.}
\begin{itemize}
\item Esta postura junto con las tres
siguientes, Triang Mukhaipāda
Paschimottānāsana, Ardha Baddha Padma
Paschimottānāsana y Paschimottānāsana,
estimulan fuertemente los órganos
abdominales. Esto nos ayuda a combatir los
males que surgen por una alimentación
inadecuada que, como dice Iyengar, a veces es
una cuestión social. Con estas cuatro
posturas rejuvenecemos los tejidos en los
órganos, glándulas y músculos de esta región.
Esto es clave para envejecer saludablemente.

Estas tres posturas también fortalecen la
columna y flexibilizan las piernas.
\item Incrementa apertura de cadera.
\end{itemize}

\emph{Mahā Mudrā.} Una postura parecida a la
descrita en el paso 4 se llama Mahā Mudrā
(Mahā = gran, Mudrā = cerradura, sello). La
pierna izquieda va flexionada de la misma
forma que en Jānu Sīrsāsāna, sólo que ahora
la espnilla sí forma un ángulo de 90º con el
muslo derecho. Se toma el dedo gordo del pie
y se estira bien la espalda, con barbilla al
pecho. En esta postura todos los órganos
entre el ano y el diaframa se tensan,
acercando el abdomen hacia la columna. Este
movimiento tiene importantes beneficios para
el sistema digestivo y las glángulas en la
región abdominal.

\subsubsection*{Triaṅg Mukhaipāda Paschimottānāsana}

Trianga = tres extremidades o partes. Se
refiere a los pies, las rodillas y los
glúteos. Mukha = cara, boca. Eka = uno.
Pāda=pie. La cara o boca tocan una pierna.
Paschimottānāsana=estiramiento de la parte
oeste (posterior) del cuerpo.

\emph{Instrucciones.}
\begin{enumerate}
\item Entrar a Daṇḍāsana.
\item Flexionar la pierna izquierda en
Vīrāsana, y mantener la derecha bien estirada
y activa con el pie flexionado.
\item Estirar ambos brazos y tomar el pie. Si
es posible, tomar con una mano la muñeca de
la otra. Otros agarres son simplemente tomar
la planta del pie, entrelazar los dedos de
las manos, tomar el dedo gordo del pie, o
alcanzar espinilla o incluso el muslo. Sea
cual sea el agarre, mantener la columna bien
estirada y el pecho abierto. Respirar
profundo.
\item En una exhalación entrar a una flexión
llevando abdomen a muslo, esternón a rodilla
y garganta a pecho.
\item Para salir primero estirar la espalda y
abrir el pecho, luego liberar los brazos y
por último las piernas. Hacer pequeños
movimientos para relajar las rodillas y hacer
el otro lado.
\end{enumerate}

\emph{Acciones.}
\begin{itemize}
\item El estiramiento de la espalda es como
en la postura anterior.
\item Cuando se alcance mucha flexibilidad,
cuidar que los codos no descansen en el piso.
\end{itemize}

\emph{Beneficios.}
\begin{itemize}
\item La flexión al frente trae los mismos
beneficios que la postura anterior.
\item Mejora condiciones en los arcos de los
pies.
\item Flexibiliza tobillos y rodillas, y
ejercita el cuadríceps.
\end{itemize}

\subsubsection*{Ardha Baddha Paschimottānāsana}

Ardha=medio. Baddha = ligado, atado.
Paschimottānāsana=estiramiento de la parte
oeste (posterior) del cuerpo.

\emph{Instrucciones.}
\begin{enumerate}
\item Entrar a Daṇḍāsana.
\item Flexionar la pierna izquierda llevando
el empeine izquierdo al muslo derecho en
Ardha Padmāsana. El talón está en contacto
con el abdomen.
\item Tomar el pie izquierdo con la mano
derecha y empujar la rodilla izquierda al
piso con la mano izquierda para preparar la
articulación. Con el tiempo es posible que la
rodilla toque el piso. Tomar algunas
respiraciones.
\item Llevar el brazo derecho por detrás de
la espalda a la altura de la cadera. En una
exhalación tomar el dedo gordo del pie
izquierdo. Llevar el hombro izquierdo hacia
atrás ayuda a lograr el agarre. Si no es
posible, regresar al paso anterior.
\item Estirar la mano derecha para tomar el
pie derecho. Si es posible, tomar el talón o
la planta del pie con la palma de la mano.
También es buena opción tomar el dedo gordo
del pie con dedos índice y medio de la mano,
o la espinilla e incluso el muslo. Estirar
bien la espalda y tomar algunas
respiraciones.
\item Inhalar, estirar por completo la
columna y en la exhalación entrar a la
flexión llevando abdomen a muslo, esternón a
rodilla y barbilla a espinilla. Ésta es la
postura completa.
\item Para salir, regresar primero el torso
estirando la columna, luego liberar las
manos, devolver el torso a la vertical y por
último liberar las piernas. Hacer pequeños
movimientos para relajar las rodillas y hacer
el otro lado.
\end{enumerate}

\emph{Acciones.}
\begin{itemize}
\item Llevar la rodilla flexionada al piso.
\item Estirar la espalda como en las posturas
anteriores.
\item Con forme se va desarrollando más
flexibilidad, acercar el pie izquierdo al
ombligo. Esto incrementa los beneficios que
se obtienen en esta zona.
\item En el punto máximo de estiramiento es
posible tomar la muñeca derecha con la mano
izquierda.
\end{itemize}

\emph{Beneficios.}
\begin{itemize}
\item La flexión hacia adelante trae los
mismos beneficios que las posturas
anteriores.
\item Se desarrolla flexibilidad en la
rodilla. Esta postura se puede tomar como una
preparación para Padmāsana.
\item Se recomienda para personas con
problemas en los hombros.
\end{itemize}

\subsubsection*{Paschimottānāsana}

Paschima = oeste, que se refiere a la parte
posterior del cuerpo desde la cabeza hasta
los talones. El este (Pūrva) es la parte
anterior del cuerpo, el norte la coronilla y
el sur las plantas de los pies y los
tobillos. Esta postura también se llama
Ugrāsana (urgā = formidable, poderoso, noble)
y Bramacharyāsana (Bramacharya = estudio
religioso, contención y celibato)

\emph{Instrucciones.}
\begin{enumerate}
\item Entrar a Daṇḍāsana.
\item Estirar los brazos y con el dedo índice
y medio de cada mano tomar el dedo gordo del
pie correspondiente manteniendo la espalda
bien estirada.
\item En una exhalación profunda flexionar
los codos y jalando con los brazos entrar a
una flexión, como en las posturas anteriores,
llevando abdomen a muslos, esternón a pecho y
frente a rodillas estirando sobre todo las
vértebras lumbares.
\item Dependiendo de la flexibilidad que se
tenga, se puede cambiar el agarre a
(enunciamos en orden ascendente de
exigencia): manos tomando talones, palmas de
las manos en plantas de los pies, una mano
toma la otra muñeca.
\item Para salir, estirar primero la espalda
como en el paso 3, y después liberar las
manos y relajar los pies.
\end{enumerate}

\emph{Acciones.} Estirar la espalda como en
las posturas anteriores.

\emph{Beneficios.}
\begin{itemize}
\item La flexión hacia adelante trae los
mismos beneficios que las posturas
anteriores.
\item En este caso, la flexión hacia adelante
trae beneficios en el corazón. Esto mejora la
circulación en todo el cuerpo.
\end{itemize}

\subsubsection*{Adho Mukha Vīrāsana}

Adho Mukha=la cara hacia abajo. Vīra=héroe,
guerrero.

\emph{Instrucciones.}
\begin{enumerate}
\item Entrar a Vīrāsana (o a Vajrāsana).
\item Llevar las palmas de las manos al piso
mantenido la columna bien estirada.
\item Exhalando, lentamente deslizar las
manos hacia en frente y hacer una flexión de
la espalda desde las vértebras lumbares.
Entrando desde Vīrāsana, llevar el torso en
medio de las piernas al piso, y entrando
desde Vajrāsana llevar el torso a los muslos.
\item Llevar primero la frente al piso y
cuando esto sea fácil, llevar la barbilla y
el esternón al piso.
\item Después de unas respiraciones, regresar
las palmas de las manos frente a las rodillas
y apoyarse en ellas para devolver la espalda
a la posición vertical.
\end{enumerate}

\emph{Acciones.}
\begin{itemize}
\item Respirar bien profundo y estirar más la
espalda desde las vértebras lumbares hasta
los brazos, tratando de avanzar un poco más
lejos en cada exhalación.
\item Devolver la mente al cuerpo.
\end{itemize}

\emph{Beneficios.}
\begin{itemize}
\item Es una postura relajante que ayuda a
combatir el cansancio causado por otras
posturas.
\item Es perfecta para practicar la atención
plena.
\item Se curva ligeramente la espalda, por lo
que sirve de contrapostura para extensiones
profundas como Uṣṭrāsana.
\end{itemize}

\subsection{Posturas de extensión}
\label{section-posturas-de-extension}

\subsubsection*{Śalabhāsana}

Śalabha = langosta, saltamontes.

\emph{Instrucciones.}
\begin{enumerate}
\item Comenzar acostado boca abajo con los
brazos extendidos a los lados del cuerpo, con
las palmas mirando hacia arriba y los
empeines sobre el piso.
\item En una exhalación levantar cabeza,
pecho, costillas, brazos y piernas del piso.
Respirar profundo.
\item Para salir, lentamente y con control
regresar a la posición inicial.
\end{enumerate}

\emph{Acciones.}
\begin{itemize}
\item Activar todo el cuerpo: piernas,
rótulas, brazos, espalda.
\item Llevar la mirada hacia arriba sin
tensar el cuello.
\item En las exhalaciones elevar más el pecho
para profundizar en el trabajo de la espalda.
No es necesario elevar más las piernas.
\end{itemize}

\emph{Beneficios.}
\begin{itemize}
\item Estira la espalda trabajando sobre todo
en la zona lumbar.
\item Mejora procesos digestivos.
\item Estimula vejiga y próstata.
\item Si se tiene exceso de arco en espalda
baja, se recomienda flexionar las rodillas.
\end{itemize}

\subsubsection*{Makarāsana}

Makara = cocodrilo. Es otra variante de
Śalabhāsana.

\emph{Instrucciones.}
\begin{enumerate}
\item Entrar a Śalabhāsana.
\item Llevar las manos a la cabeza por atrás
y entrelazar los dedos. Otra variante es
apoyando los codos en el piso y apoyando la
barbilla en las manos.
\end{enumerate}

\emph{Beneficios.} La primera variante
incrementa el calor corporal. La segunda
variante es una versión más suave que puede
ser relajante al mismo tiempo que activa la
espalda.

\subsubsection*{Bhujaṇgāsana}

Bhujaṇga = serpiente, cobra.

\emph{Instrucciones.}
\begin{enumerate}
\item Comenzar acostado boca abajo con los
brazos extendidos a los lados del cuerpo, con
las palmas mirando hacia arriba y los
empeines sobre el piso.
\item Llevar las manos a la altura del pech
con los dedos juntos.
\item Apoyar los metatarsos al piso y
empujarse ligeramente hacia el frente (esto
hace que las palmas de las manos queden a la
altura de las costillas). Ativar las piernas
con los empeines al piso.
\item Inhalar y elevar el torso del piso sólo
hasta que los codos estén en un ángulo de
90º. Así se debe practicar la postura durante
muchos meses hasta que se desarrolle
verdadera flexibilidad en la espalda.
\item Cuando esto suceda, estirar un poco más
los codos y llevar la columna a una línea
vértical perpendicular al piso sin despegar
la pelvis del piso. Respirar profundamente.
\item Cuando el paso anterior se logre,
estirar por completo los codos y llevar la
cabeza y la mirada hacia atrás. Ésta es la
postura completa.
\item Para regresar, lentamente desdoblar la
espalda, vértebra por vértebra, llevando al
final la cabeza.
\end{enumerate}

\emph{Acciones.}
\begin{itemize}
\item Es importante no forzar esta postura
para proteger las vértebras lumbares.
\item En las exhalaciones hacer extensión de
vértebras dorsales tratando de no mover las
vértebras lumbares.
\item Llevar la mirada al techo sin tensar el
cuello.
\item Mantener las piernas bien activas.
\item Juntar omóplatos.
\item En la variante flexionando las
rodillas, en las exhalaciones acercar pies a
cabeza. En la postura completa hay contacto.
\end{itemize}

\emph{Beneficios.}
\begin{itemize}
\item Revitaliza la columna y mejora
condiciones en los discos.
\item Abre el pecho.
\end{itemize}

\emph{Variante.} Una variante es flexionando
las rodillas. La acción principal aquí es
llevar pies a cabeza.

\subsubsection*{Dhanurāsana}

Dhanu = arco

\emph{Instrucciones.}
\begin{enumerate}
\item Comenzar acostado boca abajo con los
brazos extendidos a los lados del cuerpo, con
las palmas mirando hacia arriba y los
empeines sobre el piso.
\item En una exhalación flexionar las
rodillas y estirando ambos brazos tomar con
las manos los empeines, los tobillos o las
espinillas. El pecho se eleva del piso.
Respirar algunas veces.
\item En otra exhalación elevar más las
piernas. Esta acción hace que los brazos se
estiren más. (Los brazos no empujan los
pies).
\item Elevar la cabeza y llevarla tan atrás
como sea posible sin tensar el cuello.
\item Por último, si es posible, juntar
muslos, rodillas y tobillos.
\item Para salir, en una exhalación liberar
las manos y volver a la postura inicial boca
abajo.
\end{enumerate}

\emph{Acciones.}
\begin{itemize}
\item En las exhalaciones elevar más las
piernas para profundizar en la postura.
\item Tratar de arquear la espalda sobre todo
el la región lumbar para no generar tensión
en lumbares.
\item No hay ningún inconveniente en que la
respiración se acelere.
\end{itemize}

\emph{Beneficios.}
\begin{itemize}
\item Es excelente para trabajar elasticidad
en la columna.
\item Como el abdomen queda apoyado en el
piso, también estimula esta región.
\end{itemize}

\subsubsection*{Ūrdhva Mukha Śvānāsana}

Ūrdhva Mukha=boca mirando hacia arriba.
Śvāna=perro.

\emph{Instrucciones.}
\begin{enumerate}
\item Comenzar acostado boca abajo con los
brazos extendidos a los lados del cuerpo, con
las palmas mirando hacia arriba y los
empeines sobre el piso.
\item Llevar las manos a la altura del pecho
con los dedos juntos.
\item Apoyar los metatarsos al piso y
empujarse ligeramente hacia el frente (esto
hace que las palmas de las manos queden a la
altura de las costillas). Separar las piernas
al ancho de caderas o un poco menos y
mantenerlas bien activas con los empeines al
piso..
\item Inhalar y elevar el torso estirando los
brazos por completo. Al hacer este movimiento
hay que cuidarse de no hiperextender las
vértebras lumbares. La cadera y las piernas
se separan del piso y sólo queda el apoyo de
los empeines y las palmas de las manos.
\item Para salir, devolver las piernas y la
columna vértebra por vértabra al piso.
Estirar los brazos a los costados como en la
postura inicial.
\end{enumerate}

\emph{Acciones.}
\begin{itemize}
\item Activar las piernas elevando rótulas y
activar glúteos.
\item Estirar bien los brazos y la espalda
tratando de llevar la columna a una línea
completamente vertical.
\item Elevar la cabeza llevando la mirada al
techo sin tensar el cuello. La garganta va
hacia adelante.
\item Mantener las piernas y los glúteos bien
activos.
\item Llevar el pecho hacia adelante.
\end{itemize}

\emph{Beneficios.}
\begin{itemize}
\item Rejuvenece la columna. Se recomienda
para personas con espalda muy rígida.
\item Los pulmones se vuelven más elásticos.
\item Incrementa la circulación en la región
pélvica.
\end{itemize}

\subsubsection*{Chatus Pādāsana}

Chatus = cuatro, cuatro veces. Pādā = pierna.

\emph{Instrucciones.}
\begin{enumerate}
\item Comenzar acostado boca arriba en
Śavasāna.
\item Exhalando flexionar las rodillas y
acercar los talones tanto como sea posible a
la cadera.
\item Si es posible tomar los tobillos con
las manos. si esto es muy difícil, se puede
usar un cinturón haciendo un \emph{loop}
alrededor de cada espinilla que sirva como
agarre para cada mano. También es posible
dejar las palmas de las manos en el suelo con
los brazos bien estirados.
\item En una inhalación profunda elevar la
cadera y con ella la espalda. Sólo queda el
apoyo de hombros, cabeza y pies. Las manos
siguen como en el paso anterior.
\item Para salir, regresar la espalda al
piso, lento y con conciencia. Al final la
cadera. Luego estirar las piernas y volver a
Śavasāna.
\end{enumerate}

\emph{Acciones.}
\begin{itemize}
\item En las exhalaciones elevar más la
cadera. Cuidar de no forzar las lumbares.
\item Mantener las espinillas totalmente
verticales, sin que las rodillas se acerquen
entre sí.
\item Relajar cuello y cintura escapular.
\end{itemize}

\emph{Beneficios.}
\begin{itemize}
\item Incrementa la flexibilidad en la
columna. Es un estiramiento muy profundo que
puede ser realizado antes de una postura
clímax como Uṣṭrāsana o Ūrdhva Dhanurāsana.
\item Elimina tensión emocional acumulada en
el pecho.
\item Abre el pecho y cintura escapular.
\item Incrementa flexibilidad en pelvis.
\end{itemize}

\subsubsection*{Uṣṭrāsana}

Uṣṭra = camello

\emph{Instrucciones.}
\begin{enumerate}
\item Comenzar en Vajrāsana.
\item Apoyar las manos en frente de las
piernas para separar las rodillas y los pies
al ancho de caderas. Hincarse.
\item Llevar las manos a la cadera con los
pulgares en dirección de los glúteos y los
demás dedos hacia las ingles.
\item Inhalar y estirar la columna llevando
coronilla al cielo, y en la exhalación
comenzar a extender la espalda lentamente
llevando la cabeza hacia atrás.
\item Cuando se haya llegado al punto máximo
de extensión tomar algunas respiraciones.
\item Llevar mano derecha a talón derecho y
mano izquierda a talón y izquierdo. Si es
posible, apoyar la palma de la mano en la
planta del pie. Ésta es la postura completa,
donde se deben realizar las acciones
mencionadas a continuación.
\item Para salir, regresar una mano a la
cadera y luego la otra. Lentamente devolver
la columna a la posición vertical, vértebra
por vértebra.
\item Por último flexionar rodillas y
regresar a Vajrāsana. Es recomentable hacer
una contra postura después de esta como Adho
Mukha Vīrāsana.
\end{enumerate}

\emph{Acciones.}
\begin{itemize}
\item En cada exhalación tratar de llevar la
cadera hacia adelante para que los muslos
estén verticales.
\item Estirar profundamente la región dorsal
de la espalda y el cóccix.
\item Activar muslos y glúteos.
\end{itemize}

\emph{Beneficios.}
\begin{itemize}
\item Estimula fuertemente la flexibilidad de
la columna.
\item Abre el pecho y flexibiliza el cóccix.
\item Se recomienda a personas con espalda
cifótica.
\end{itemize}

\subsubsection*{Ūrdvha Dhanurāsana}

Ūrdvha =arriba. Dhanu = arco. También llamada
Cakrāsana (Cakra = rueda).

\emph{Instrucciones.}
\begin{enumerate}
\item Comenzar acostado boca arriba en
Śavasāna.
\item Flexionar las rodillas y acercar los
talones a la cadera. Llevar las palmas de las
manos abajo de los hombros con los dedos
hacia los pies. Los codos miran al techo. La
distancia entre las manos no debe ser mayor
que el ancho de hombros.
\item En una exhalción profunda elevar la
cadera del piso extendiendo fuertemente la
columna. Los únicos apoyos con las palmas de
las manos y las plantas de los pies. Ésta es
la postura completa y aquí se deben realizar
las acciones.
\item Para salir, en una exalación regresar
lentamente la cadera y la espalda al piso.
Puede resultar cómodo apoyarse en los
metatarsos.
\end{enumerate}

\emph{Acciones.}
\begin{itemize}
\item Si al principio no es posible entrar a
esta postura, tener paciencia y hacer
posturas más suaves como Chatus Pādāsana.
\item Al principio puede ser que los talones
no estén en el piso. Tratar de que sí haya
contacto.
\item Estirar bien los codos y llevar la
cadera más hacia arriba para profundizar en
la extensión.
\item Si es posible, caminar con las manos
hacia los pies.
\end{itemize}

\emph{Beneficios.}
\begin{itemize}
\item Todo el cuerpo se llena de energía con
esta extensión de columna. Nos mantiene en un
estado alerta y lleno de vida.
\item Aumenta la flexibilidad de las muñecas
y de los brazos.
\item La sangre fluye a la cabeza, mejorando
la circulación de todo el cuerpo.
\item Se estiran las ingles.
\item Se liberan emociones acumuladas en el
pecho y en la región pélvica.
\item Nos da una gran sensación de ligereza.
\end{itemize}

\subsubsection*{Pūrvottānāsana}

Pūrva = este. Se refiere a la parte anterior
del cuerpo, desde la frente hasta los dedos
de los pies.

\emph{Instrucciones.}
\begin{enumerate}
\item Entrar a Daṇḍāsana.
\item Inhalar profundo y en la exhalación
elevar cadera y pecho manteniendo las piernas
bien estiradas. Si es posible llevar las
plantas de los pies al piso sin doblar
rodillas.
\item Para regresar, flexionar los codos y
las rodillas para llevar lentamente el cuerpo
al piso. Relajar las piernas.
\end{enumerate}

\emph{Acciones.}
\begin{itemize}
\item Activar todo el cuerpo de pies a cabeza
y estirarlo profundamente.
\item No doblar las rodillas aunque las
plantas de los pies no estén totalmente en el
piso. Con el tiempo y la práctica esto se
volverá posible.
\item Elevar la cadera al techo.
\item Abrir el pecho juntando omóplatos atrás
de la espalda.
\item Mantener los brazos bien estirados.
\item Relajar la cabeza.
\end{itemize}

\emph{Beneficios.}
\begin{itemize}
\item Se fortalecen los tobillos y las
muñecas.
\item Mejora las articulaciones de los
hombros.
\item Expande el pecho.
\item Es excelente para equilibrar el trabajo
realizado en las flexiones hacia adelante.
\end{itemize}

\subsection{Posturas para fortalecer abdomen}
\label{section-posturas-abdomen}

\subsubsection*{Paripūrṇa Nāvasāsana}

Paripūrṇa = completo, entero. Nāva = bote,
barco, embarcación.

\emph{Instrucciones.}
\begin{enumerate}
\item Entrar a Daṇḍāsana..
\item Exhalar y llevar el tronco ligeramente
hacia atrás y levantar las piernas totalmente
activas y estiradas. Pies en punta. Todo el
peso descansa en los isquiones y glúeos, nada
en la columna o espalda. En la postura
completa en ángulo entre piernas y torso es
de 60-65º. Los pies están por arriba de la
cabeza.
\item Levantar los brazos estirados paralelos
al piso. Las palmas mirándose mutuamente,
cinco dedos de la mano juntos.
\item Sostener la postura 30 segundos e ir
aumentando hasta llegar a un minuto. Los
efectos se sienten al pasar los 20 segundos.
\item Exhalar y devolver piernas y brazos al
piso. Lentamente llevar la espalda al piso y
descansar boca arriba.
\end{enumerate}

\emph{Acciones.}
\begin{itemize}
\item Activar abdomen.
\item Estirar espalda.
\item Llevar brazos paralelos al piso.
\item Elevar las piernas sin doblar rodillas
ni perder cualquiera de las acciones
anteriores.
\end{itemize}

\emph{Beneficios.}
\begin{itemize}
\item Ayuda a aliviar hinchazón en el
estómago. Como Pávana Muktāsana, sirve para
aliviar gases en el estómago.
\item Ayuda a reforzar espalda; muy
recomendado para mujeres embarazadas, pues
ayuda para cargar al bebé. Refuerza espalda
baja.
\item En contraste con Ardha Nāvāsana, es un
trabajo fuerte en intestinos.
\end{itemize}

\subsubsection*{Ardha Nāvasāsana}

Ardha = medio, mitad. Nāva = bote, barco,
embarcación.

\emph{Instrucciones.}
\begin{enumerate}
\item Entrar a Daṇḍāsana.
\item Entrelazar dedos atrás de la cabeza
justo arriba del cuello.
\item Exhalar y llevar el tronco ligeramente
hacia atrás y levantar las piernas
totalemente activas y estiradas. Pies en
punta. Todo el peso descansa en los isquiones
y glúeos, nada en la columna o espalda. El
ángulo entre piernas y torso es de 30-35º.
Los pies están a la altura de la cabeza.
\item Sostener la postura 30 segundos e ir
aumentando hasta llegar a un minuto. Los
efectos se sienten al pasar los 20 segundos.
\item Si se sostiene el aire en esta postura
el trabajo se intensifica en el esómago y no
en los músculas abdominales. Se puede
practicar respirando profundamente o haciendo
inhalaciones y exhalaciones intercaladas con
retenciones.
\end{enumerate}

\emph{Acciones.}
\begin{itemize}
\item Activar abdomen.
\item Estirar espalda.
\item Llevar brazos paralelos al piso.
\end{itemize}

\emph{Beneficios.}
\begin{itemize}
\item Ayuda a aliviar hinchazón en el
estómago. Como Pávana Muktāsana, sirve para
aliviar gases en el estómago.
\item Ayuda a reforzar espalda; muy
recomendado para mujeres embarazadas, pues
ayuda para cargar al bebé. Refuerza espalda
baja.
\item En contraste con Paripūrṇa Nāvāsana, es
un trabajo fuerte en hígado, vesícula y bazo.
\end{itemize}

\subsubsection*{Ūrdvha Prasarita Pādāsana}

Ūrdvha = arriba. Prasarita = extendido,
separado. Pādā = pie.

\emph{Instrucciones.}
\begin{enumerate}
\item Comenzar acostado boca arriba.
\item Opcional: colocar los brazos abajo del
cuerpo con las palmas mirando hacia abajo y
los dedos bien separados. Esto facilita la
postura.
\item Activar ambas piernas bien estiradas
con los pies en punta.
\item Elevar las piernas sólo 30º y sostener
algunas respiraciones.
\item Elevar las piernas hasta los 60º y
sostener algunas respiraciones.
\item Elevar las piernas hasta los 90º y
estirar ambos brazos hacia atrás. Ésta es la
postura completa.
\item Para salir regresar primero los brazos
y luego bajar las piernas bien estiradas,
lentamente y con control.
\end{enumerate}

\emph{Acciones.}
\begin{itemize}
\item Activar abdomen y piernas.
\item Estirar bien los brazos y columna.
\item Al principio puede ser difícil llevar
las dos piernas al ángulo de 90º. Esto tiene
que ver con falta de flexibilidad en
isquiotibiales. Con el tiempo se desarrolla
esta flexibilidad. Hacer la postura de forma
relajada y sin forzar el estiramiento
facilitará el proceso.
\item Un excelente ejercicio es entrar y
salir de esta postura varias veces. Esto
ayuda a fortalecer el abdomen y las piernas.
\item Otro ejercicio recomendable es entrar a
esta postura (sin detenerse en el ascenso de
las piernas) y continuar el movimiento hasta
llegar a Halāsana. Luego regresar las piernas
estiradas hasta llegar casi al suelo, sin
tocarlo, y volver a subir. Repetir alrededor
de diez veces para preparar el abdomen para
otras posturas como Salaāmba Sarvāngāsana I.
\end{itemize}

\emph{Beneficios.} Fortalece fuertemente el
abdomen y las piernas, y también la espalda
baja (cuadrado lumbar). También trabaja sobre
órganos abdominales.

\subsubsection*{Jaṭhara Parivartanāsana}

Jaṭhara = abdomen. Parivṛtta = girado,
torcido.

\emph{Instrucciones.}
\begin{enumerate}
\item Acostado con las piernas activas, pies
flexionados y brazos como avión estirados a
los lados, palmas mirando hacia abajo.
\item Flexionar rodillas y acercar al pecho.
\item Rotar las rodillas y la cadera hacia un
lado, la cabeza mira en dirección contraria.
\item Hacer ambos lados.
\end{enumerate}

\emph{Acciones.}
\begin{itemize}
\item Jaṭhara Parivartanāsana A es con las
rodillas dobladas, con las instrucciones que
hemos dado.
\item Jaṭhara Parivartanāsana B es con las
piernas estiradas. Para entrar, comenzar
acostado con piernas y brazos como avión,
como en la variante A. Levantar las piernas a
90º, y girar hacia un lado. La cabeza mira en
dirección contraria. Hacer ambos lados.
\end{itemize}

\emph{Posturas preparatorias.} Pávana
Muktāsana, Gato-Vaca (Bidālāsana o
Mārjārīāsana), Balāsana.

\emph{Beneficios.} Elimina rigidez en hombros
y columna, estimula sistema digestivo, libera
tensión en espalda.

\subsection{Posturas de equilibrio}
\label{section-posturas-de-equilibrio}

\subsubsection*{Vasiṣṭhāsana}

Vasiṣṭha fue un sabio de la antigüedad. Se le
atribuyen muchos himnos Védicos, entre los
que destaca el séptimo Maṇḍala del
\d{R}g Veda.
Era representativo de la dignidad y el poder
brahmánicos, y es uno de los siete sabios
identificados con las estrellas del Gran Oso.
La rivalidad entre él y el sabio real
Viśvāmitra, un kṣatriya (hombre de la raza de
los guerreos) que, por su piedad y ascetismo
ascendió al estado brahmánico, es motivo de
varias leyendas.

\emph{Instrucciones.}
\begin{enumerate}
\item Entrar a Adho Mukha Śvānāsana.
\item Apoyar la mano derecha con los dedos
bien abiertos justo abajo del centro del
pecho.
\item Guiado por un movimiento de apertura en
el pecho (el hombro izquierdo hacia el
techo), girar todo el cuerpo hasta que el
arco externo del pie derecho y el interno del
pie izquierdo estén apoyados en el piso.
\item Si es posible apoyar un pie sobre el
otro de forma que los únicos apoyos sean el
arco externo del pie derecho y la mano
derecha.
\item Elevar el brazo izquierdo vertical al
cielo en una misma línea con el brazo derecho
y llevar la mirada hacia el techo.
\item Para salir regresar primero el brazo
izquierdo al piso para tener estabilidad y
luego ambos pies al piso. Con control volver
a Adho Mukha Śvānāsana y hacer el otro lado.
\end{enumerate}

\emph{Acciones.}
\begin{itemize}
\item Busca que se forme una línea recta en
el costado izquierdo desde la pierna hasta la
axila.
\item Mantener la columna y el torso bien
rectos.
\item Abrir el pecho usando la apertura de la
mano que está arriba.
\item Llevar la mirada más hacia arriba sin
tensar cuello.
\item Si no es posible llevar un pie sobre el
otro se pueden dejar ambos al piso.
\item Si no es posible elevar la mano que no
está en el piso se puede dejar descansando
sobre el costado de arriba.
\end{itemize}

\subsubsection*{Chaturaṅga daṇḍāsana}

Chatur = cuatro. Ariga = extremidad, parte.
Daṇḍa = bastón, vara. Los cuatro soportes son
las manos y los pies.

\emph{Instrucciones.}
\begin{enumerate}
\item Comenzar acostado boca abajo.
\item Llevar las palmas de las manos a los
costados del pecho. Apoyar metatarsos de los
pies al piso.
\item Activar todo el cuerpo, sobre todo las
piernas y el abdomen, y en una exhalación
separarse del piso apoyándose sólo en las
manos y en los metatarsos.
\item Respirar profundo y salir en una
exhalación para volver boca abajo sobre el
piso.
\end{enumerate}

\emph{Acciones.}
\begin{itemize}
\item Activar piernas y abdomen.
\item Mirada al frente.
\item Codos tocando el pecho.
\item Espalda bien estirada.
\item Encontrar la relajación aunque la
postura sea muy activa.
\end{itemize}

\emph{Beneficios.} Fortalece el abdomen, las
muñecas y los brazos. Es una postura de
transición en los saludos al sol.

\subsubsection*{Bakāsana}

Baka = cuervo.

\emph{Instrucciones.}
\begin{enumerate}
\item Comenzar en Varjrāsana.
\item Apoyar los metatarsos y las manos en el
piso.
\item Separar las rodillas más que el ancho
de homrbos y las caderas.
\item Dejar que el peso del cuerpo se incline
hacia adelante hasta doblar los codos. Apoyar
las rodillas en la parte posterior de los
brazos cerca de las axilas.
\item Levantar un pie y buscar el equilibrio,
y luego elevar el otro pie.
\item Estirar los brazos y mantener el
equilibrio.
\item Para salir, volver a flexionar los
codos y llevar el peso atrás para volver a
Varjrāsana.
\end{enumerate}

\emph{Acciones.}
\begin{itemize}
\item Llevar mirada al frente y estirar la
espalda.
\item Pies en punta lejos de cadera.
\item Conforme se vaya perfeccionando la
postura se acerca más el apoyo de las
rodillas hacia las axilas.
\item Una entrada alternativa es desde
Sīrsāsana .
\item Una salida alternativa es dar un brinco
hacia atrás hasta Chaturaṅga daṇḍāsana.
\end{itemize}

\emph{Beneficios.}
\begin{itemize}
\item Foralece los brazos y el abdomen.
\item Nos sirve para practicar una postura en
equilibrio como preparación para las
inversiones.
\item Estimula los órganos abdominales.
\end{itemize}

\emph{Variantes.} Una variante con torsión es
Pārśva Bakāsana, con piernas dobladas, con
una pierna estirada o con piernas estiradas.

Otra variente importante Eka Pāda Bakāsana,
cuando tenemos un sólo apoyo de rodilla a
codo (hay que trabajar ambos lados). Como
mencionamos al inicio, hay una entrada a
Bakāsana desde Sīrsāsana II, aunque esta
última no está incluida en este manual.

\subsection{Posturas invertidas}
\label{section-posturas-invertidas}

\subsubsection*{Adho Mukha Śvānāsana}

Adho Mukha = boca mirando hacia abajo. Śvāna
= perro.

\emph{Variantes.} Rodillas flexionadas para
estirar espalda. Al estirar rodillas se
pierde el estiramiento de espalda.

\emph{Instrucciones.}
\begin{enumerate}
\item Comenzar acostado boca abajo.
\item Flexionar los codos y colocar las
palmas de las manos sobre el piso a la altura
del pecho. Los dedos índices están paralelos
al lado más largo del mat y los demás dedos
bien separados.
\item En una exhalación elevar el torso y
llevar la cadera hacia atrás, estirando bien
los codos. La cabeza cuelga y el cuello está
relajado. Respirar profundo y realizar las
acciones.
\item Para salir, flexionar las rodillas y
volver a la posición boca abajo.
\end{enumerate}

\emph{Acciones.}
\begin{enumerate}
\item Al principio, cuando estiramos las rodillas la
espalda se redondea. En esta etapa es
preferible flexionar las piernas para que la
espalda esté lo más recta posible.
Concentrarse en estirar espalda baja y
permitir que los isquiotibiales se relajen y
flexibilicen para acercar talones al piso. El
proceso puede durar años.
\item Hacer rotación externa con ambos brazos
(esto lo puede hacer el instructor) para
abrir bien el pecho.
\item Estirar conscientemente la espalda
baja.
\item Isquiones al techo.
\item Llevar la mente a la espalda: debe
estar recta como una tabla.
\item Cuando ya se tiene flexibilidad en
isquitoibiales, profundizar acercando talones
al piso.
\item Cuando ya se tiene flexibilidad en
todas las acciones anteriores, acercar la
frente al piso hasta que la cabeza se apoye
en el piso sin arquear la espalda (Iyengar).
\end{enumerate}

\emph{Beneficios.}
\begin{itemize}
\item Quita la fatiga.
\item Por su acción en isquiotibiales, es
buena para corredores.
\item Fortalece tobillos y piernas.
\item Abre el pecho y trabaja cintura
escapular (aperturas de pecho pueden ser
buenas posturas y ejercicios preparatorios).
\item Disminuye el ritmo cardiaco. Es
relajante.
\item Al ser una inversión, nos ayuda a
enfrentar miedos y puede ser una postura
preparatoria para Sīrsāsana I.
\item También por ser una inversión, mejora
la circulación en la cabeza.
\item Las personas con presión alta, a
diferencia de otras inversiones, sí pueden
practicar esta postura.
\end{itemize}

\subsubsection*{Adho Mukha Vrksāsana}

Adho Mukha = boca mirando hacia abajo. Vrksa
= árbol.

\emph{Instrucciones.}
\begin{enumerate}
\item Comenzar en Tadāsana.
\item Flexionar el cuerpo hacia adelante
hasta apoyar las manos en el piso a distancia
de hombros. No doblar los codos.
\item En una exhalación llevar los pies hasta
que estén verticales.
\end{enumerate}

\emph{Acciones.}
\begin{itemize}
\item Separar hombros de orejas.
\item Estirar espalda.
\item Activar abdomen.
\item Pies en punta.
\item Abrir el pecho.
\item Al principio resultará indispensable
hacer la postura contra una pared. Las
primeras veces subir con una pierna y luego
la otra. Luego se puede experimentar subir
doblando ambas rodillas.
\item Practicar con frecuencia esta postura
ayuda a perfeccionarla.
\end{itemize}

\emph{Beneficios.} Armoniza el cuerpo.
Fortalece brazos y muñecas, y ayuda a abrir
el pecho.

\subsubsection*{Pīncha Mayūrāsana}

Pīncha = barbilla o pluma. Mayūra =
pavorreal. Cuando es temporada de lluvias,
los pavorreales bailan, y elevan sus plumas
formando un abanico: eso es lo que representa
esta postura.

\emph{Instrucciones.}
\begin{enumerate}
\item Comenzar hincado. Flexionar hacia
adelante y apoyar palmas y antebrazos en el
piso al ancho de hombros. Los antebrazos
deben estar paralelos.
\item Elevar ambas piernas a la vertical.
Pies en punta.
\end{enumerate}

\emph{Acciones.}
\begin{itemize}
\item Estirar bien el cuello y mirar la pared
de atrás. En una variante inicial la barbilla
va ligeramente hacia el pecho, incluso
mirando los pies. Más adelante se pueden
explorar extensiones en esta postrura y la
mirada va hacia las manos.
\item Piernas bien activas.
\item Hombros lejos de orejas.
\item Estirar completamente la columna.
\item No permitir que los brazos dejen de
estar paralelos. Al principio es muy
conveniente usar un bloque en donde se
encajen las comisuras de los dedos índice y
pulgar de cada mano. También se puede usar un
cinturó en los antebrazos para evitar que
éstos se separen.
\item Al principio resultará indispensable
hacer la postura apoyándose en una pared. Se
puede explorar subir elevanto una pierna y
luego la otra.
\item Practicar con frecuencia esta postura
ayuda a perfeccionarla.
\end{itemize}

\emph{Beneficios.} Fortalece cintura
escapular, espalda, columna y abdomen.

\subsubsection*{Sālamba Sarvāngāsana I}

Ālamba = prop, soporte. Sa = acompañado de.
Sālamba = con un soporte. Sarva = todo,
completo. Anga = extremidad, miembro, cuerpo.
Es la postura del cuerpo completo en un
soporte.

\emph{Instrucciones.}
\begin{enumerate}
\item Comenzar en Śavāsana.
\item Activar piernas, pies en punta y
estirar los brazos a los costados con las
palmas mirando hacia el piso. Respirar bien
profundo.
\item Flexionar rodillas y llevar los muslos
al abdomen acercando lo más posible. Respirar
profundamente.
\item Elevar suavemente la cadera del piso y
flexionar los codos para que la cadera se
apoye en las manos. Respirar profundamente y
explorar la comodidad en esta postura.
\item Exhalando elevar la cadera para que la
espalda esté completamente vertical sin
estirar las rodillas. La barbilla toca el
pecho. Reajustar el soporte de las manos para
que esté a la altura de las vértebras
dorsales. Respirar muy profundamente.
\item Exhalando estirar ambas piernas hasta
que estén completamente verticales. Las
rodillas bien activas y los pies en punta.
\item Mantener la postura el tiempo que sea
posible respirando profundo y permitiendo que
el cuerpo se relaje profundamente.
\end{enumerate}

\emph{Acciones.}
\begin{itemize}
\item Estirar por completo el cuerpo desde la
columna hasta las puntas de los pies. Buscar
una línea completamente vertical. Es
conveniente reajustar muchas veces el soporte
de las manos.
\item Activar abdomen.
\item Respirar bien profundo y hacer
conscientes los procesos que suceden en el
cuerpo y la mente. Puede ser sumamente
relajante.
\item Es posible entrar a la postura con
ambas piernas estiradas llevándolas primero
hacia Halāsana (aunque no hace falta que los
pies toquen el piso, basta que las piernas
estiradas estén parlelas al piso) y luego a
la vertical.
\end{itemize}

\emph{Beneficios.} Es de las mayores
bendiciones que nos fueron transmitidas por
los sabios antiguos. Sarvāngāsana es la madre
de todas las posturas. Estimula el sistema
endócrino al estimular las glándulas tiroides
y paratiroides, que están situadas en la
región del cuello. Es una inversión que no
exige tanta práctica como otras en cuanto a
balance y fuerza, pero aún así trae los
beneficios principales de las inversiones:
ayuda la circulación, el ritmo cardiaco, el
sistema nervioso y la respiración. Citando a Iyengar: ``no hay
exageración en decir que una persona que
practica Sarvāngāsana regularmente sentirá
vigor y fuerza, y será feliz y tendrá
confianza en sí mismo''.

\emph{Contraindicaciones.} Las personas que
sufren de presión alta no deben practicar
esta postura a menos de que puedan sostener
Halāsana por al menos tres minutos.

\subsubsection*{Eka Pāda Sarvāngāsana}

Eka = uno. Pāda = pie.

\emph{Instrucciones.}
\begin{enumerate}
\item Entrar a Sarvāngāsana.
\item Bajar una pierna al piso en posición de
Halāsana mientras la otra se mantiene
estirada y completamente vertical.
\item Elevar el pie para volver a
Sarvāngāsana y hacer el otro lado.
\end{enumerate}

\emph{Acciones.} Al principio será difícil
que el pie de abajo toque el piso. Ser
paciente y permitir que los músculos
encuentren la flexibilidad que necesitan a
través de la relajación y la respiración.
Después de algunos segundos, en las
exhalaciones, se puede ir acercando cada vez
más el pie al piso.

\emph{Beneficios.} Se estimulan los riñones y
se fortalecen las piernas.

\subsubsection*{Pārśv Eka Pāda Sarvāngāsana}

Pārśva = lado,costado. Eka = uno. Pāda = pie.

\emph{Instrucciones.}
\begin{enumerate}
\item Entrar en Sarvāngāsana y hacer Eka Pāda
Sarvāngāsana con ambas piernas y volver a la
vertical.
\item Exhalar y bajar la pierna derecha hacia
el lado derecho completamente estirada de
forma que ambas piernas estén en un plano
perpendicular al piso. El pie de arriba está
completamente vertical. Respirar
profundamente.
\item Regresar la pierna derecha a la
vertical y hacer la pierna izquierda.
\end{enumerate}

\emph{Acciones.} Como en Eka Pāda
Sarvāngāsana, al principio será difícil que
el pie de abajo toque el piso. Ser paciente y
permitir que los músculos encuentren la
flexibilidad que necesitan a través de la
relajación y la respiración. Después de
algunos segundos, en las exhalaciones, se
puede ir acercando cada vez más el pie al
piso.

\emph{Beneficios.} Se estimulan los riñones y
se fortalecen las piernas, ayuda con procesos
digestivos.

\emph{Otras variantes de Sālamba
Sarvāngāsana.} Baddha Koṇāsana en
Sarvāngāsana, Ūrdhva Padmāsana en
Sarvāngāsana y Pinḍāsana en Sarvāngāsana.

\subsubsection*{Halāsana}

Hala = arado.

\emph{Instrucciones.}
\begin{enumerate}
\item Entrar a Ūrdvha Prasārita Pādāsana.
\item Llevar lentamente y con conciencia las
piernas atras de la cabeza hasta que toquen
el piso. Si es posible apoyar los empeines.
La espalda debe estar estirada y
perpendicular al piso.
\item Para salir devolver muy lentamente las
piernas al piso (estiradas si es posible),
sintiendo el contacto de cada vértebra con el
piso.
\end{enumerate}

\emph{Acciones.}
\begin{itemize}
\item Al principio puede ser complicado que
los pies toquen el piso. Es recomentable usar
un soporte (silla) para irse acostumbrando a
la postura.
\item Cerrar los ojos y centrar la mirada en
la punta de la nariz.
\item Estirar bien la espalda perpendicular
al piso.
\item Llevar la manos detrás de la cabeza y
entrelazar los dedos de las manos para abrir
el pecho. Intentar que las palmas de las
manos en el entrelazado miren hacia el techo
para profundizar en la apertura.
\item Es posible entrar a la postura desde
Sarvāngāsana:
\begin{itemize}
\item Flexionar ligeramente la espalda y
llevar simultáneamente los brazos y las
piernas por detrás de la cabeza. Estirar bien
las piernas.
\item Devolver los brazos al soporte de
Sarvāngāsana a y luego estirarlos
completamente sobre el piso. Si es posible
entrelazar los dedos de las manos para abrir
el pecho.
\end{itemize}
\end{itemize}

\emph{Beneficios.} Comparte muchos beneficios
con Sarvāngāsana, en particular estimula
fuertemente el sistema endócrino al estimular
las glándulas tiroides y paratiroides, que
están situadas en la región del cuello. Ayuda
a desarrollar confianza en uno mismo al ser
una inversión que se puede hacer con soporte,
obteniendo gran estabilidad.

\emph{Contraindicaciones.} Las personas que
sufren de presión alta deben practicar esta
postura con particular atención. La
respiración no debe ser agitada y no se debe
sentir presión en la cabeza. No dudar en usar
soporte.

\subsubsection*{Karṇapīdāsana}

Karṇa = oreja. Pīda = dolor, incomodidad o
presión.

\emph{Instrucciones.}
\begin{enumerate}
\item Después de hacer Halāsana por algún
tiempo, flexionar las rodillas y acercarlas a
las orejas. Tanto las rodillas como los
empeines están en contacto con el piso.
\item Para salir regresar a Halāsana o
levantar las rodillas flexionadas para
después llevarlas al piso y terminar en
Śavāsana.
\end{enumerate}

\emph{Acciones.}
\begin{itemize}
\item Al principio puede ser complicado que
las rodillas toquen las orejas sin despegar
los pies del piso. Ser paciente y respirar
profundo ayudará a aumentar la flexibilidad.
\item Se puede intentar tomar los talones con
las manos, y también es bueno estirar los
brazos hacia atrás entrelazando manos para
abrir más el pecho.
\item Cerrar los ojos y centrar la mirada en
la punta de la nariz.
\item Estirar bien la espalda perpendicular
al piso.
\end{itemize}

\emph{Beneficios.} Comparte muchos beneficios
con Sarvāngāsana, en particular estimula
fuertemente el sistema endócrino al estimular
las glándulas tiroides y paratiroides, que
están situadas en la región del cuello. Ayuda
a la circulación, descansa la espalda, el
tronco y las piernas.

\emph{Contraindicaciones.} Las personas que
sufren de presión alta deben practicar esta
postura con particular atención. La
respiración no debe ser agitada y no se debe
sentir presión en la cabeza.

\subsubsection*{Supta Koṇāsana}

Supta = acostado. Koṇa = ángulo.

\emph{Instrucciones.}
\begin{enumerate}
\item Después de hacer Karṇapīdāsana por
algún tiempo, estirar la piernas abriendo el
compás tanto como sea posible. Se apoyan los
metatarsos en el piso.
\item Para salir regresar a Halāsana o
levantar las rodillas flexionadas para
después llevarlas al piso y terminar en
Śavāsana.
\end{enumerate}

\emph{Acciones.}
\begin{itemize}
\item Al principio puede ser complicado que
los pies toquen el piso. Ser paciente y
respirar profundo ayudará a aumentar la
flexibilidad.
\item Se puede intentar tomar los dedos
gordos de los pies con las manos. Esto puede
ayudar a profundizar, pero la espalda se debe
mantener vertical.
\item Cerrar los ojos y centrar la mirada en
la punta de la nariz.
\item Estirar bien la espalda perpendicular
al piso.
\end{itemize}

\emph{Beneficios.} Además de los beneficios
de las posturas anteriores, fortalece las
piernas.

\emph{Contraindicaciones.} Las personas que
sufren de presión alta deben practicar esta
postura con particular atención. La
respiración no debe ser agitada y no se debe
sentir presión en la cabeza.

\subsubsection*{Pārśva Halāsana}

Pārśva = lado, costado. Hala = arado.

\emph{Instrucciones.}
\begin{enumerate}
\item Después de hacer Supta Koṇāsana por
algún tiempo, volver a Halāsana.
\item Apoyar las manos en la espalda como
sporte.
\item Llevar ambas piernas juntas y con las
rodillas bien estiradas hacia el lado
derecho. Tras algunas respiraciones y hacer
el otro lado.
\end{enumerate}

\emph{Acciones.} Estirar bien las piernas y
mantener la espalda tan vertical como sea
posible.

\emph{Beneficios.}
\begin{itemize}
\item Es una compactación muy profunda de los
órganos abdominales. Ayuda con procesos de
eliminación en el sistema digestivo
(intestinos).
\item Aumenta la flexibilidad de la columna.
\item Comparte beneficios con las posturas
anteriores como estimular glándulas tiroides
y paratiroides.
\end{itemize}

\emph{Contraindicaciones.} Las personas que
sufren de presión alta deben practicar esta
postura con particular atención. La
respiración no debe ser agitada y no se debe
sentir presión en la cabeza.

\subsubsection*{Sālamba Sīrsāsana I}

Ālamba = prop, soporte. Sa = acompañado de.
Sālamba = con un soporte. Sīrsa = cabeza.

\emph{Instrucciones.}
\begin{enumerate}
\item Comenzar hincado y apoyar los
antebrazos en el piso tomando con cada mano
el codo opuesto. Ésta es la medida de los
codos, no deben moverse a partir de este
punto. Entrelazar los dedos de las manos
liberando sólo los dedos meñiques.
\item Elevar las caderas y estirar rodillas.
\item Respirando lentamente y con control
levantar los pies del piso.
\item Estirar la espalda.
\item Elevar los muslos a la vertical sin
desdoblar las rodillas. Los pies miran a la
pared de atrás y las espinillas están
paralelas al piso.
\item Elevar las espinillas para llegar a la
postura completa.
\item Para salir hacer exactamente los mismos
pasos en el orden invertido.
\end{enumerate}

\emph{Acciones.}
\begin{enumerate}
\item Estirar completamente la espalda. La
columna, los isquiotibiales y los talones
están alineados verticalmente. La garganta,
la barilla y el esternón también están
alineados verticalmente.
\item Los bordes externos e internos de las
manos también están alineados verticalmente
para que el soporte de la cabeza sea
correcto.
\item Separar hombros de orejas.
\item Meter cóccix ligeramente, como en
Tādāsana.
\item Buscar que el cuerpo esté completamente
vertical.
\item Relajar el cuello, no debe haber
tensión.
\item No permitir que se separen los codos.
\item Respirar profundamente y relajar la
mente.
\item El peso debe estar en la cabeza y no en
los brazos, que sólo sirven para mantener el
equilibrio.
\item Pies en punta o pies en flex (se pueden
explorar ambas opciones. Iyengar indica pies
en punta.)
\item Se puede explorar otro entrelzado de
dedos cuando se tenga buen equilibrio: las
palmas de las manos y las muñecas juntas. Es
posible con este entrelazado no sea necesario
liberar el dedo meñique que queda abajo.
\end{enumerate}

\emph{Observaciones.}
\begin{enumerate}
\item Al principio resulta muy útil usar una
pared para desarrollar el balance necesario.
También es recomendable asistencia de un
maestro o amigo con experiencia.
\item Cuando ya se tiene buena estabilidad,
se puede entrar a la postura con las piernas
bien estiradas pasando por Ūrdvha Dandāsana.
También se puede hacer esto para salir.
\item El tiempo que se puede pasar en esta
postura depende de cada quién. Se puede pasar
hasta quince minutos sin bajar haciendo las
variantes que presentaremos a continuación.
Puede ser que al principio sea mucho menos,
apenas unos segundos.
\item Es recomendable sentirse muy cómodo en
las dos posturas anteriores, Sarvāngāsana y
Halāsana antes de comenzar a practicar
Sīrsāsana I.
\item Se recomienda hacer Sālamba
Sarvāngāsana después de Sīrsāsana I para
evitar volverse irritable
\end{enumerate}

\emph{Beneficios.}
\begin{enumerate}
\item Ayuda a enfrentar miedos.
\item Desarrolla el equilibrio en sentido
físico y más allá. Quien la practica se
vuelve balanceado y resistente al dolor y al
placer, la pérdida y la ganancia, la pena y
la fama, y la derrota y la victoria.
\item La cabeza es el centro de las
cualidades sátvicas, el contexto de los tres
gunas (satva, rajas y tamas). Esta postura
rejuvenece las células en el cerebro de forma
que los pensamientos y las ideas son más
claros e inmediatos. Mejora la vitalidad.
\item Amplía los horizontes del espíritu.
\item Mejora circulación, presión,
respiración, sistema inmunológico, sistema
neriovoso.
\item Estimula las grlándulas pineal y
pituitaria.
\item Ayuda a combatir insomnio, letargo,
mala memoria, resfriados, tos, palpitaciones.
\item Junto con Sarvangāsana ayuda a combatir
constipación.
\item Incremente la energía de manera global.
\end{enumerate}

\emph{Contraindicaciones.} No se recomienda
si hay presión alta o baja.

\subsubsection*{Eka Pāda Sīrsāsana}

Eka = uno. Pāda = pie.

\emph{Instrucciones.}
\begin{enumerate}
\item Entrar a Sīrsāsana I.
\item Bajar una pierna al piso mientras la
otra se mantiene estirada y completamente
vertical.
\item Elevar el pie para volver a Sīrsāsana y
hacer el otro lado.
\end{enumerate}

\emph{Acciones.}
\begin{itemize}
\item Al principio será difícil que el pie de
abajo toque el piso. Ser paciente y permitir
que los músculos encuentren la flexibilidad
que necesitan a través de la relajación y la
respiración. Después de algunos segundos, en
las exhalaciones, se puede ir acercando cada
vez más el pie al piso.
\item Es posible que el cuello sienta mucha
presión. Para evitar esto, mantener ambas
piernas bien activas, la de arriba
completamente vertical y la de abajo sólo
hasta el punto el que sea cómodo.
\item Entre cada pierna respirar tantas veces
como sean necesarias para encontrar la calma.
\item No doblar las rodillas, puesto esto
causa pérdida de equilibrio.
\end{itemize}

\emph{Beneficios.} Se estimulan los riñones y
se fortalecen las piernas. También se
fortalecen el cuello y el abdomen.

\subsubsection*{Pārśva Eka Pāda Sīrsāsana}

Pārśva = lado,costado. Eka = uno. Pāda = pie.

\emph{Instrucciones.}
\begin{enumerate}
\item Entrar en Sīrsāsana I.
\item Exhalar y bajar la pierna derecha hacia
el lado derecho completamente estirada de
forma que ambas piernas estén en un plano
perpendicular al piso. El pie de arriba está
completamente vertical. Respirar
profundamente.
\item Regresar la pierna derecha a la
vertical y hacer la pierna izquierda.
\end{enumerate}

\emph{Acciones.}
\begin{itemize}
\item Como en Eka Pāda Sīrsāsana, al
principio será difícil que el pie de abajo
toque el piso. Ser paciente y permitir que
los músculos encuentren la flexibilidad que
necesitan a través de la relajación y la
respiración. Después de algunos segundos, en
las exhalaciones, se puede ir acercando cada
vez más el pie al piso.
\item Mantener ambas piernas bien estiradas.
\item Hacer consciente la flexión en la ingle
de la pierna que está flexionada.
\end{itemize}

\emph{Beneficios.} Fortalece el cuello, el
abdomen, los muslos y la columna. Estimula
riñones e intestinos.

\subsubsection*{Pārśva Sīrsāsana}

Pārśva = lado,costado. Pāda = pie.

\emph{Instrucciones.}
\begin{enumerate}
\item Entrar a Entrar a Sīrsāsana I.
\item Manteniendo la columna totalmente
estirada, girar el cuerpo hacia un lado hasta
llegar, si es posible, a los 90º.
\item Exhalar, regresar a Sīrsāsana y hacer
el otro lado.
\end{enumerate}

\emph{Acciones.}
\begin{itemize}
\item Una y otra vez asegurarse de que la
columna esté recta.
\item Mantener los pies en flex.
\item Al principio resulta difícil hacer el
movimiento, pero con el tiempo se desarrolla
confianza para profundizar.
\end{itemize}

\emph{Beneficios.} Además de los beneficios
de Sīrsāsana I, aumenta la flexibilidad de la
espalda.

\subsubsection*{Parivṛtta Eka Pāda Sīrsāsana}

Parivṛtta = girado, torcido. Eka = uno. Pāda
= pie.

\emph{Instrucciones.}
\begin{enumerate}
\item Entrar a Sīrsāsana I y hacer Eka Pāda
Sīrsāsana con la pierna derecha.
\item Llevar la pierna al lado izquierdo
haciendo una torsión.
\item Salir de la torsión, regresar a
Sīrsāsana I y hacer el otro lado.
\end{enumerate}

\emph{Acciones.}
\begin{itemize}
\item Las piernas completamente estiradas y
bien activas.
\item En las exhalaciones, acercar la pierna
de abajo al piso y profundizar en la torsión.
\item Mantener la espalda bien estirada.
\item Alejar hombros de orejas y relajar el
cuello.
\item Iyengar indica que la pierna de arriba
tiene que ir hacia atrás, de manera que la
línea vértical sea la bisectriz del ángulo en
las piernas.
\end{itemize}

\emph{Acciones.} Al principio será difícil
que el pie de abajo toque el piso. Ser
paciente y permitir que los músculos
encuentren la flexibilidad que necesitan a
través de la relajación y la respiración.
Después de algunos segundos, en las
exhalaciones, se puede ir acercando cada vez
más el pie al piso.

\emph{Beneficios.} Además de los beneficios
de Sīrsāsana I, se estimulan los riñones y se
fortalecen las piernas.

\subsubsection*{Baddha Koṇāsana Sīrsāsana}

Baddha = ligado, atado. Koṇa = ángulo.

\emph{Instrucciones.}
\begin{enumerate}
\item Desde Sīrsāsana I flexionar rodillas y
juntar plantas de los pies.
\item Respirar profundamente y volver a
Sīrsāsana I.
\end{enumerate}

\emph{Acciones.} Acercar talones a pelvis sin
tensar el cuello, manteniendo hombros lejos
de orejas y columna bien estirada.

\emph{Beneficios.} Además de los beneficios
de Sīrsāsana I, aumenta la flexibilidad en
cadera.

\subsubsection*{Ūrdhva Koṇāsana en Sīrsāsana}

Ūrdva = arriba. Koṇa = ángulo.

\emph{Instrucciones.}
\begin{enumerate}
\item Desde Sīrsāsana I, al exhalar abrir el
compás manteniendo ambas piernas bien
estiradas. Los pies van en dirección del
piso.
\item Después de llegar al máximo punto,
flexionar ligeramente la cadera sin generar
tensión en cabeza. Si es posible llevar los
metatarsos al piso.
\item Regresar al paso anterior y finalmente
a Sīrsāsana I.
\end{enumerate}

\emph{Acciones.}
\begin{itemize}
\item También es posible salir de Sīrsāsana I
a través de esta postura cuando los pies
tocan el piso.
\item Cuidar que el cuello no se tense por
intentar llevar las piernas al piso. La
flexibilidad se desarrolla con el tiempo y la
postura se debe lograr con facilidad.
\end{itemize}

\emph{Beneficios.} Además de los beneficios
de Sīrsāsana I, aumenta la flexibilidad en
caderas y fortalece isquiotibiales.

\subsubsection*{Ūrdva Padmāsana en Sīrsāsana}

Padma = loto

\emph{Instrucciones.}
\begin{enumerate}
\item Desde Sīrsāsana I flexionar una rodilla
y llevar el pie al muslo opuesto. Respirar
profundo.
\item Si es posible hacerlo sin demasiado
esfuerzo, cuidando de no tensar el cuello,
llevar el otro pie al costado opuesto para
entrar en Padmāsana.
\item Regresar a Sīrsāsana I y hacer el otro
lado.
\end{enumerate}

\emph{Acciones.}
\begin{itemize}
\item Acercar las rodillas entre sí y estirar
los muslos para profundizar en la postura.
\item Puede ser sumamente relajante, respirar
bien profundo y disfrutar.
\item También es posible hacer torsión en
esta postura: Pārśva Ūrdhva Padmāsana en
Sīrsāsana. En esta postura se esimulan aún
más los órganos abdominales y mejora la
circulación en esta área.
\end{itemize}

\emph{Beneficios.}
\begin{itemize}
\item Se estira la pelvis, toda la espalda y
las costillas.
\item Se expande el pecho.
\item Mejora circulación en la pelvis.
\end{itemize}

\subsubsection*{Pinḍāsana en Sīrsāsana}

Pinḍā = embrión.

\emph{Instrucciones.}
\begin{enumerate}
\item Desde Ūrdhva Padmāsana en Sīrsāsana
simplemente flexionar la cadera y llevar las
rodillas a las axilas.
\item Regresar a Ūrdhva Padmāsana en
Sīrsāsana, luego a Sīrsāsana y hacer el otro
lado.
\end{enumerate}

\emph{Acciones.} Acercar más las rodillas a
las axilas respirando profundo y sin
sobrepasar los límites del cuerpo.

\emph{Beneficios.} Las mismas que la postura
anterior, pero con una contracción aún mayor
en la parte abdominal del cuerpo.

\subsubsection*{Ūrdhva Daṇḍāsana}

Ūrdvha = arriba. Daṇḍa = apoyo, bastón, vara.

\emph{Instrucciones.}
\begin{enumerate}
\item Desde Sīrsāsana I llevar ambas piernas
juntas y estiradas hacia el piso hasta que
estén paralelas al piso.
\item Regresar a Sīrsāsana I.
\end{enumerate}

\emph{Acciones.}
\begin{itemize}
\item Al principio será difícil hacer que
ambas piernas estén paralelas al piso. Ser
paciente y permitir que los músculos
encuentren la flexibilidad que necesitan a
través de la relajación y la respiración.
Después de algunos segundos, en las
exhalaciones, se puede ir profundizando más
sin ir más allá de la horizontal y sin
generar tensión en el cuello.
\item Como en Daṇḍāsana, estirar
completamentente la espalda y activar bien
las piernas.
\item Relajar bien el cuello alejando hombros
de orejas.
\item La espalda no debe estar perpendicular
al piso: se inclinará un poco hacia atrás
para compensar el peso de las piernas.
\item También es posible salir (y entrar) a
Sīrsāsana a través de esta postura. Se
recomienda hacerlo cuando ya se tiene muy
buena estabilidad.
\item Se puede sostener hasta por un minuto
incluso si es sólo postura de transición para
entrar o salir de Sīrsāsana.
\end{itemize}

\emph{Beneficios.} Se estimulan los riñones y
se fortalecen las piernas.

\medskip\noindent\emph{Variaciones con silla}

\subsubsection*{Halāsana con silla}

(Ver Halāsana.) Hala = arado.

\emph{Variantes.} Brazos estirados adelante
(también se pueden entrelazar dedos de las
manos para abrir pecho) o brazos estirados
hacia atrás.

\emph{Instrucciones.}
\begin{enumerate}
\item Colocar una cobija (con pocos dobleces,
no queremos mucha altura) en la mitad del
mat.
\item Colocar una silla en un extremo del mat
con el asiento mirando hacia la cobija.
\item Acostarse boca arriba con los hombros y
las últimas vértebras sobre la cobija pero la
cabeza sobre el mat.
\item Doblar las rodillas y poner plantas de
los pies al piso.
\item Estirar bien los brazos y poner palmas
de las manos al piso.
\item Lentamente y con control elevar la
cadera del piso y llevar las piernas hacia
atrás hasta que toquen la silla.
\item Si es necesario, quizás con ayuda del
instructor, ajustar la distancia de la silla
para que los pies estén en el asiento.
\item Entrelazar dedos de las manos atrás de
la espalda para profundizar en la apertura
del pecho. También se pueden llevar los
brazos atrás.
\end{enumerate}

\emph{Acciones.}
\begin{itemize}
\item Elevar la cadera y estirar bien la
espalda. Intentar que la columna esté en una
línea recta perpendicular al piso.
\item Activar y estirar bien las piernas.
\item Respirar profundo.
\end{itemize}

\emph{Beneficios.}
\begin{itemize}
\item Como en Halāsana, se estimulan
fuertemente las glándulas tiroides y
paratiroides, que están situadas en la región
del cuello.
\item Ayuda a desarrollar fuerza en las
piernas y espalda.
\item Permite la circulación a la cabeza.
\item Abre el pecho.
\item Esta postura es muy cómoda y puede ser
una excelente manera para comenzar a trabajar
Halāsana y Sarvāngāsana de forma segura.
Recomendada para personas con muy poca
flexibilidad, y para quienes necesitan
desarrollar seguridad con las inversiones.
\end{itemize}

\subsubsection*{Sarvāngāsana con silla}

(Ver Sālamba Sarvāngāsana I.) Sarva = todo,
completo. Anga = extremidad, miembro, cuerpo.
Es la postura del cuerpo completo en un
soporte.

\emph{Instrucciones.}
\begin{enumerate}
\item Colocar con la parte enfrente a la
mitad del mat. Una cobija va, sin muchos
dobleces como en Halāsana con silla en el
piso cerca de la silla. Otra cobija en el
asiento de la silla cubriendo los bordes del
mismo.
\item Sentarse con los isquiones sobre la
cobila que está en el asiento y las piernas
dobladas sobre el respaldo. El borde del
respaldo va a en la parte de atrás de las
rodillas.
\item Tomar fuertemente el respaldo de la
silla y llevar el peso hacia atrás hasta que
la espalda baja está sobre el borde de
enfrente del asiento y la cabeza cuelgue.
\item Deslizar el cuerpo hasta que la cabeza
toque el piso y después los hombros. Las
manos también deben ir deslizando cambiando
el agarre a las patas de enfrente de la
silla.
\item Al terminar el deslizamiento la espalda
alta debe estar en la misma posición que en
Halāsana con silla: los hombros y las
primeras vértebras dorsales sobre la cobija y
la cabeza en el piso. El sacro es apoya en la
cobija sobre la silla.
\item En este punto las piernas aún están
apoyadas en el respaldo. Estirarlas y hacer
algunas respiraciones.
\item Por último, estirar las piernas hacia
arriba: ésta es la postura completa.
\item Para salir, primero devolver las
piernas sobre el respaldo de la silla.
\item Lentamente deslizar la cabeza y los
hombros hacia atrás lejos de la silla. Cuando
sea necesario apoyar las plantas de los pies
sobre el respaldo de la silla y después sobre
el asiento.
\item Eventualmente la cadera tocará el piso
y las pantorrillas estarán sobre el asiento.
Ésta es una postura relajante y se puede
sostener por mucho tiempo.
\item Continuar deslizándose hacia atrás
hasta que todo el cuerpo esté sobre el piso.
\end{enumerate}

\emph{Acciones.}
\begin{enumerate}
\item Estirar bien las piernas hacia arriba
en una línea perpendicular al piso. Los pies
en punta.
\item Estirar la espalda pero exagerar la
curvatura lumbar.
\item Relajar homrbos y cuello.
\item Si es posible tomar las piernas
traseras de la silla.
\item Respirar profundo y disfrutar el
soporte de la silla: la postura no es tan
intensa.
\end{enumerate}

\emph{Beneficios.}
\begin{itemize}
\item Como en Sarvāngāsana, Halāsana y
Halāsana con silla y, se estimulan
fuertemente las glándulas tiroides y
paratiroides, que están situadas en la región
del cuello.
\item Fortalece las piernas y abre pecho.
\item Permite la circulación a la cabeza.
\item Sumamente relajante, se puede sostener
por mucho tiempo.
\end{itemize}

\subsubsection*{Sīrsāsana con silla}

(Ver Sālamba Sīrsāsana I.) Sīrsa = cabeza.

\emph{Instrucciones.}
\begin{enumerate}
\item Colocar dos sillas a un costado de la
pared, una frente a la otra, a distancia de
unos treinta o cuarenta centímetros. En cada
asiento colocar una cobija sin muchos
dobleces.
\item Comenzar hincado frente a las dos silla
mirando la pared.
\item Inclinarse hacia en frente levantando
la cadera y lentamente apoyando los hombros
en cada asiento de la silla. La cabeza en el
espacio entre las sillas.
\item Tomar los bordes de las sillas con las
manos.
\item Acercar lo más que se pueda los pies a
la pared, aún apoyados sobre el piso.
\item Flexionar una rodilla y luego la otra.
\item Estirar bien la espalda y pegarla a la
pared con las rodillas dobladas.
\item Estirar las piernas hacia arriba y
recargarlas en la pared.
\item Para salir hacer exactamente los mismos
pasos en orden invertido.
\end{enumerate}

\emph{Acciones.}
\begin{itemize}
\item Estirar bien las piernas hacia arriba.
Los pies pueden ir flexionados o en punta.
\item Relajar los hombros en el soporte.
Asegurarse de que no haya nada de tensión
aquí y si es necesario salir de la postura y
reajustarla.
\item Ayudarse del apoyo de las manos para
encontrar estabilidad y relajarse.
\end{itemize}

\emph{Beneficios.} Es un primer acercamiento
a la sensación de Sīrsāsana. Nos sirve para
desarrollar equilibrio y confianza en las
inversiones.

\subsection{Saludo al Sol}
\label{section-saludo-al-sol}

\subsubsection*{Saludo al Sol A}

\emph{Secuencia.}
\begin{enumerate}
\item Tadāsana
\item Tadāsana con Ūrdhva Hastāsana
\item Uttānāsana
\item Ardha Uttānāsana
\item Brinco hacia atrás
\item Plancha (Phalakāsana)
\item Chaturaṅga daṇḍāsana
\item Urdhva Mukha Śvānāsana
\item Chaturaṅga daṇḍāsana
\item Adho Mukha Śvānāsana
\item Talones arriba y salto adelante
\item Uttānāsana
\item Ardha Uttānāsana
\item Tadāsana con Ūrdhva Hastāsana
\item Tadāsana
\end{enumerate}

\subsubsection*{Saludo al Sol B}

\emph{Secuencia.}
\begin{enumerate}
\item Tadāsana
\item Utkatasana B
\item Uttānāsana
\item Ardha Uttānāsana
\item Brinco hacia atrás
\item Plancha (Phalakāsana)
\item Chaturaṅga Daṇḍāsana
\item Urdhva Mukha Śvānāsana
\item Chaturaṅga Daṇḍāsana
\item Adho Mukha Śvānāsana
\item Vīrabhadrāsana I con pierna derecha
\item Chaturaṅga Daṇḍāsana
\item Adho Mukha Śvānāsana
\item Chaturaṅga Daṇḍāsana
\item Urdhva Mukha Śvānāsana
\item Adho Mukha Śvānāsana
\item Vīrabhadrāsana I con pierna izquierda
\item Chaturaṅga daṇḍāsana
\item Urdhva Mukha Śvānāsana
\item Adho Mukha Śvānāsana. Cinco
respiraciones largas
\item Talones arriba y salto adelante
\item Uttānāsana
\item Utkatāsana B
\item Tadāsana
\end{enumerate}

\subsubsection*{Saludo al Sol Clásico}

\emph{Secuencia.}
\begin{enumerate}
\item Tadāsana
\item Tadāsana con Ūrdhva Hastāsana y
extensión
\item Uttānāsana
\item Ānjaneyāsana
\item Plancha (Phalakāsana)
\item Bhujaṇgāsana
\item Ānjaneyāsana
\item Uttānāsana
\item Tadāsana con Ūrdhva Hastāsana y
extensión
\item Samasthiti (Tadāsana con Namaskarāsana)
\item Tadāsana
\end{enumerate}

\subsection{Posturas Restaurativas}
\label{section-posturas-restaurativas}

\subsubsection*{Viparīta Karaṇī}

Viparīta = invertido. Karaṇī= prática. Es la
postura del gesto (Mudrā) invertido. Se
considera una āsana y también un Mudrā.

\emph{Variantes.} Brazos estirados a los
lados con las palmas mirando hacia abajo,
brazos hacia atrás tomándose de los codos, o
brazos estirados hacia atrás.

\emph{Instrucciones.}
\begin{enumerate}
\item Lo más seguro es que necesitemos un
soporte para realizar esta postura. Cada
persona puede necesitar distintos soportes.
Se puede usar una, dos o hasta tres cobijas,
o un bolster. Colocar el soporte en la pared
sobre el mat.
\item Con la cadera sobre el soporte y las
rodillas flexionadas sobre un costado (en
posición fetal), acercar los isquiones a la
pared.
\item Llevar la cadera completamente sobre el
soporta como en Pávana Muktāsana y estirar
las piernas completamente hacia el techo con
los pies flexionados.
\item Colocar los brazos en cualquiera de las
variantes. La postura es más profunda
estirando los brazos hacia atrás.
\end{enumerate}

\emph{Acciones.}
\begin{itemize}
\item Estirar complatamente la espalda
separando vértebra por vértebra.
\item Relajar el cuello.
\item Activar y estirar ambas piernas.
\end{itemize}

\emph{Beneficios.}
\begin{itemize}
\item Esta postura es una inversión, por lo
que comparte varios beneficios otras
invertsiones (ver Sīrsāsana) como:
\begin{itemize}
\item Permite la circulación desde los pies
hacia la cabeza. Calma la mente.
\item Reactiva sistema linfático.
\item Mejora sistema endócrino (glándulas).
\item Combate la fatiga y nos da energía
\end{itemize}
\item Al mismo tiempo es una postura
restaurativa, por lo que se puede lograr una
relajación muy profunda y activación del
sistema nervioso parasimpático.
\item Desarrolla flexibilidad en
isquiotibiales y espalda.
\end{itemize}

\emph{Contraindicaciones.} Evitar durante la
menstruacioón.

\subsubsection*{Setu Bhanda Sarvāngāsana}

Setu = puente. Setu Bhanda = construcción de
un puente.

\emph{Instrucciones.}
\begin{enumerate}
\item Colocar el mat pegado a la pared, un
bloque en altura alta pegado a la pared y
dejar otro bloque cerca.
\item Comenzar acostado boca arriba y
flexionar las rodillas llevando las plantas
de los pies al piso.
\item En una exhalación elevar la cadera y
depositar un bloque debajo del sacro (nunca
sobre vértebras lumbares) en la altura que
sea más cómoda. Respirar profundo y
relajarse. Conforme pasa el tiempo la columna
alcanza mayor flexibilidad y puede ser
posible subir el bloque a la altura media o
alta. Ésta es una postura en sí misma que
incluimos abajo como una variante.
\item Estirar ambas piernas llevando los pies
sobre el bloque con las plantas de los pies
pegadas a la pared. Dejar los brazos a los
costados con las palmas mirando hacia a
arriba o entrelazar los dedos de las manos
para abrir más el pecho.
\item Para salir, primero devolver los pies
al piso, elevar ligeramente la cadera y poner
el bloque a un lado. Devolver la columna al
piso lentamente, vértebra por vértebra, con
conciencia. Por último estirar las piernas al
frente.
\end{enumerate}

\emph{Acciones.}
\begin{itemize}
\item Respirar profundo y relajar el cuerpo.
Dedicarse sobre todo a relajar la mente y
observarse a uno mismo.
\item Elevar el esternón y rotar externamente
los brazos para profundizar en la apertura de
pecho.
\item No forzar vértebras lumbares.
\item Activar bien las piernas.
\end{itemize}

\emph{Beneficios.}
\begin{itemize}
\item Es profundamente relajante. Nos ayuda a
recuperar energía.
\item Incremente la flexibilidad en la
columna.
\item Abre el pecho y relaja el cuello.
\item Estimula el sistema nervioso.
\end{itemize}

\emph{Variante con bolster.} En esta variante
se usa un bolster en lugar del bloque va en
el sacro. La espalda baja y media se
recuestan sobre el bolster, y los hombros y
la cintura escapular están bien apoyados en
el piso como en Sarvāngāsana. Es una postura
sumamente amigable y relajante. Perfecta para
principiantes y personas de avanzada edad.

\emph{Postura del medio puente.} Esta
variante se distingue de la postura completa
por mantener las rodillas flexionadas y los
pies al piso. Como dijimos arriba, puede
comenzar con un bloque en la altura más baja
y conforme pase el tiempo ir incrementando la
altura. En la opción más profunda se coloca
un bloque en altura alta sobre otro en altura
baja.

El bloque siempre debe estar en contacto con
el sacro y nunca con las vértebras lumbares,
que se podrían lastimar si no se coloca
adecuadamente. Esta variante se puede
sostener por mucho tiempo, por lo que es muy
buena opción para terminar una práctica.
También se puede incorporar en una clase
suave agregando el ejercicio de estirar una
pierna y luego la otra.

\emph{Postura completa.} La postura completa
se puede considerar una variante
Sarvāngāsana. Los brazos y la cintura
escapular están exactamente como en esa
postura, pero las piernas van estiradas hacia
el piso, con las plantas de los pies bien
apoyadas. Es una postura exigente.

\subsubsection*{Supta Baddha Koṇāsana}

Supta = acostado. Baddha = ligado, atado.
Koṇa = ángulo. (Ver Baddha Koṇāsana)

\emph{Variantes.} Cinturón y bolster,
cinturón sin bolster, bolster y brazos atrás,
o bolster y brazos a los lados.

\emph{Instrucciones.} Mostramos las
instrucciones para entrar a la variante con
cinturón:
\begin{enumerate}
\item Colocar el bolster sobre el mat, a lo
largo, cerca de uno de los bordes.
\item Doblar una cobija como acordeón a lo
largo en cuatro pliegues. Colocarla sobre el
bolster (a lo largo, no transversalmente) a
partir de la mitad de manera que casi toda la
cobija quede sobre el piso.
\begin{itemize}
\item Dependiendo de la flexibilidad en
espalda baja o condiciones como lordosis
puede ser conveniente usar dos cobijas. En
este caso doblar otra cobija como la primera,
colocándola por debajo del bolster de tal
forma que uno de los extremos sirva como
asiento.
\end{itemize}
\item Sentarse en Baddha Koṇāsana a distancia
de una mano del bolster.
\item Colocar el cinturón de la siguiente
manera:
\begin{itemize}
\item Por atrás de la espalda pasa a la
altura del sacro. \textbf{No colocar sobre
vértebras lumbares.}
\item Por adelante pasa por abajo de los
bordes externos de los pies.
\item El engarzado debe ser entre el pie y el
muslo, donde el cinturón no hace contacto con
el cuerpo. Es indistinto de qué lado.
\item Ajustar tanto como sea posible sin
forzar el cuerpo. Esto nos ayuda a
desarrollar la flexibilidad en caderas sin
hacer esfuerzo.
\end{itemize}
\item Apoyar las manos en el piso y vértebra
por vértebra recostarse sobre el bolster. La
cobija debe estar a la altura de las
vértebras dorsales. Es probable que en un
primer intento la cobija no esté a la altura
correcta. Sin pensarlo dos veces, salir de la
postura para corregir esto las veces que sean
necesarias (es probable que se esté en la
postura durante mucho tiempo y vale la pena
entrar correctamente).
\item Cuando la espalda esté en posición
correcta, tomar con las manos el extremo
libre de la cobija y enrollarlo comenzando
\emph{hacia abajo}. Hacer el rollo lo que se
suficiente para que sirva como soporte para
la cabeza.
\item Relajarse profundamente. Respirar y
traer la mente a la postura. Observar la
mente, las sensaciones físicas y las
emociones. Percibir sin juzgar. Respirar más.
\item Después de unos minutos en la postura
se abre el pecho naturalmente y es posible
llevar los brazos atrás: elevar los brazos
frente al pecho, tomarse de los codos y
llevar detrás de la cabeza.
\item Cuando haya transcurrido la mitad del
tiempo que se quiere estar en la postura
regresar los brazos arriba del pecho para
cambiar el entrelazado y llevar nuevamente
hacia atrás.
\item Para salir regresar los brazos a los
costados y apoyar las palmas al piso. Elevar
lentamente la espalda vértebra por vértebra
hasta hestar sentado.
\item Aflojar el cinturón y luego deshacer el
engarzado.
\end{enumerate}

\emph{Acciones.}
\begin{itemize}
\item Estirar bien la espalda y relajar el
cuello.
\item Permitir conscientemente que se abra el
pecho.
\item Llevar la mente a las caderas para que
se flexibilicen.
\item Ligeramente meter el cóccix para
profundizar en la flexión de caderas.
\item Respirar bien profundo y realizar la
postura conscientemente.
\end{itemize}

\emph{Beneficios.}
\begin{itemize}
\item Es muy poderosa como relajación. El
sistema nervioso parasimpático se activa, nos
reconocemos en un espacio seguro lejos de
todo peligro.
\item Esto permite al cuerpo puede realizar
funciones restaurativas naturalmente
\item Se desarrolla la flexibilidad de las
caderas y la espalda a través de la
relajación y no del trabajo.
\item Es exelente como primera postura en una
práctica para cambiar de entorno.
\end{itemize}

\subsubsection*{Supta Vīrāsana}

Supta=acostado. Vīra=héroe, guerrero. (Ver
Vīrāsana)

\emph{Instrucciones.}
\begin{enumerate}
\item Realizar exactamente la misma
configuración de bolster y cobija(s) que en
Supta Baddha Koṇāsana. De hecho, las
indicaciones para entrar y estirar los brazos
en la postura son exactamente las mismas
salvo que ahora se entra desde Vīrāsana.
\end{enumerate}

\emph{Acciones.} Mismas que en Supta Baddha
Koṇāsana salvo que ahora al meter el cóccix
se profundiza en profundizar en el
estiramiento de cuadríceps (muslos). Para
profundizar en este estiramiento se trata
acercar las rodillas. Recordar que en esta
variante restaurativa es priotiaria la
relajación y la respiración que desarrollar
la flexibilidad.

\emph{Beneficios.} Comparte los beneficios de
Supta Baddha Koṇāsana, salvo que en este caso
no hay trabajo en caderas sino en cuadríceps
(muslos) y rodillas.

\emph{Contraindicaciones.} Para una persona
que necesita restaurar la movilidad de
rodillas y empeines ésta es una buena
alternativa a Supta Vīrāsana. Si hay dolor
intenso en estas partes del cuerpo mejor no
hacer la postura.

\subsubsection*{Supta Sukhāsana}

Supta=acostado. Sukha=plenitud, amplitud.
(Ver Sukhāsana)

\emph{Instrucciones.}
\begin{enumerate}
\item Realizar exactamente la misma
configuración de bolster y cobija(s) que en
Supta Baddha Koṇāsana. De hecho, las
indicaciones para entrar y estirar los brazos
en la postura son exactamente las mismas
salvo que ahora se entra desde Sukhāsana.
\end{enumerate}

\emph{Acciones.} Mismas que en Supta Baddha
Koṇāsana. Ésta es una vairante muy amable que
puede ser una alternativa si hay dolor en
rodillas, empeines o cadera.

\emph{Beneficios.} Sumamente relajante. (Ver
Supta Baddha Koṇāsana)

\subsubsection*{Dwi Pāda Viparīta Daṇḍāsana}

Dwi = dos. Pā dā = pierna, pie. Viparīta =
invertido. Daṇḍa = bastón, vara.

\emph{Instrucciones.} Mostramos cómo entrar a
esta postura en su variante con silla.
\begin{enumerate}
\item Colocar la silla con el respaldo a
distancia un poco menor que un metro de la
pared sobre el piso. No se debe usar mat en
esta configuración, pues la silla debe
resbalar hacia atrás. Colocar un bloque en la
altura máxima pegado a la pared frente a la
silla y otro bloque unos 40 centímetros atrás
de la silla también en altura máxima.
\item Sentarse en la silla con el pecho hacia
el respaldo.
\item Tomar el respaldo con ambas manos,
inhalar y estirar bien la espalda haciendo
tracción de columna.
\item En la exhalación llevar lentamente la
espalda hacia atrás asegurando el agarre de
las manos. Llegar a la máxima extensión y
dejar que la cabeza caiga hacia atrás. Tomar
la silla de los lados del respaldo.
\item Llevar los pies sobre el bloque con las
plantas de pegadas a la pared. Estirar las
piernas de manera que la silla resbale hacia
atrás la medida justa.
\item Seguramente el bloque no estará
exactamente abajo de la cabeza y habrá que
ajustarlo. A veces es posible hacerlo uno
mismo y a veces es necesario pedir ayuda a un
compañero. La coronilla debe estar en
contacto con el bloque.
\item Tomar los codos frente al pecho y
llevar los brazos hacia atrás. Esta puede ser
la postura final, que es sumamente profunda
en cuanto apertura de pecho. Otra opción es
estirar ambos brazos y tomar las patas
traseras (las que están más cerca de la
pared) de la silla. Para que el agarre de las
manos nos ayude a abrir el pecho debemos
dirigir las palmas externamente de forma que
los pulgares apunten al piso. Si esto no es
posible se pueden dejar las manos descansando
sobre el piso.
\item Para salir primero soltar las manos de
las patas traseras. Podemos ayudarnos de las
patas delanteras, del asiento y eventualmente
del respaldo de la silla para reincorporar la
espalda, vértebra por vértebra, hasta
terminar con la columna vertical.
\item Antes de salir de la silla se
recomienda hacer la contrapostura: apoyar un
antebrazo sobre el borde superior del
respaldo de la silla, y el otro antebrazo
sobre el primero. Descansar la frente sobre
los antebrazos y encorvar la espalda.
Respirar bien profundo el tiempo que se
necesario y luego salir lentamente de la
silla.
\end{enumerate}

\emph{Acciones.}
\begin{itemize}
\item La postura debe ser sumamente cómoda.
Al principio parece desafiante: esto es
simple apariencia. Una vez que nos
familiarizamos con la configuración es una
postura sencilla y sumamente relajante.
\item Se puede estar bastante tiempo en esta
postura. Esto ayuda a la columna y el pecho a
flexibilizarse sin que nosotros hagamos
esfuerzo, de manera que podemos profundizar
simplemente relajando el cuerpo y la mente.
\item Activar bien las piernas y mantener los
pies en flex pegados, tanto como sea posible,
a la pared.
\end{itemize}

\emph{Beneficios.}
\begin{itemize}
\item La apertura de pecho y el estiramiento
de columna son sumamente profundos.
\item La sangre fluye a la cabeza, mejorando
la circulación en todo el cuerpo.
\item Se fortalecen las piernas
\item La mente y el cuerpo entran en un
estado de relajación profunda.
\end{itemize}

En la postura completa, los brazos están en
posición de Sīrsāsana I y las plantas de los
pies están apoyadas en el piso. Se trata de
una postura muy profunda que no incluimos en
este manual.

\subsubsection*{Śavāsana}

Śava, Mṛta = cadáver. (Esta postura también
es llamada Mṛtāsana).

\emph{Instrucciones.}
\begin{enumerate}
\item Acostarse boca arriba con las piernas
extendidas y los brazos a los costados.
\item Las palmas mirando hacia arriba y el
pecho ligeramente abierto.
\item Los brazos no tocan el resto del
cuerpo.
\item Los pies caen a los lados y los talones
se tocan.
\end{enumerate}

\emph{Acciones.}
\begin{itemize}
\item Concentrarse. Traer la mente al cuerpo,
al presente. Percibir atentamente
pensamientos, sensaciones físicas, emociones,
sonidos, estímulos de luz, etc., \emph{sin
juzgar}, sin involucrarse, sino apenas
permitiendo todo lo que acontece.
\item Una y otra vez volver a la
concentración.
\item Respirar profundamente, con conciencia,
desde diafragma. En cada exhalación relajar
el cuerpo físicamente, rostro, mandíbula,
hombros, brazos, torso, cadera, piernas,
pies. En cada exhalación traer de vuelta la
mente a la práctica.
\end{itemize}

\emph{Beneficios.}
\begin{itemize}
\item Calma la mente y elimina la fatiga que
pueden generar otras posturas. Es una
relajación profunda.
\item Permite el flujo del Prāṇa a través de
los Nāḍīs.
\end{itemize}

\bibliography{my}
\bibliographystyle{amsalpha}

\end{document}
